\exercisesection{Complements of Linear Algebra}

\begin{enumerate}
\def\labelenumi{\arabic{enumi})}
\tightlist
\item
  Let
\end{enumerate}

\[
\mathbf {M} \rightarrow \mathbf {N} \xrightarrow {t} \mathbf {P} \rightarrow \mathbf {Q}
\]

\[
\begin{array}{c c c c c} \downarrow u & & \downarrow v & & \downarrow w \\ \downarrow & & \downarrow & & \downarrow s \end{array}
\]

\[
\mathrm{M} ^ {\prime} \rightarrow \mathrm{N} ^ {\prime} \xrightarrow {t ^ {\prime}} \mathrm{P} ^ {\prime} \rightarrow \mathrm{Q} ^ {\prime}
\]

a commutative diagram of A-modules where the rows are exact; suppose that \(u\) is surjective and \(s\) injective. Show that:

\[
\operatorname{Ker} (w) = t (\operatorname{Ker} (v)) \quad \operatorname{Im} (v) = t ^ {\prime - 1} (\operatorname{Im} (w)).
\]

\begin{enumerate}
\def\labelenumi{\arabic{enumi})}
\setcounter{enumi}{1}
\tightlist
\item
  Let \(u: \mathbf{M} \to \mathbf{N}, v: \mathbf{N} \to \mathbf{P}\) two homomorphisms of A-modules. Show that there exists an exact sequence:
\end{enumerate}

\[
0 \rightarrow \operatorname{Ker} (u) \rightarrow \operatorname{Ker} (v \circ u) \rightarrow \operatorname{Ker} (v) \rightarrow \operatorname{Coker} (u) \rightarrow \operatorname{Coker} (v \circ u) \rightarrow \operatorname{Coker} v \rightarrow 0.
\]

\begin{enumerate}
\def\labelenumi{\arabic{enumi})}
\setcounter{enumi}{2}
\tightlist
\item
  In the diagram of A-modules:
\end{enumerate}

\[
\begin{array}{c} \mathrm{P} \xleftarrow {} \mathrm{M} \\ \Big \downarrow \quad \mathrm{R} ^ {\swarrow} \Big \downarrow \\ \mathrm{N} \xleftarrow {} \mathrm{Q} \end{array}
\]

Does the commutativity of the two triangles extracted from the diagram necessarily imply that of the diagram?

\begin{enumerate}
\def\labelenumi{\arabic{enumi})}
\setcounter{enumi}{3}
\item
  Let \(k\) a field, \(B = k[T]\) , let A be the subring \(k[T^2, T^3] \subset k[T]\) . Show that the A-module \(B\) is torsion-free, but is not a flat A-module.
\item
  Let \(k\) a field, \(\mathbf{C} = k[\mathbf{X}, \mathbf{Y}]\) ; consider the subrings \(\mathbf{A} = k[\mathbf{X}^4, \mathbf{Y}^4]\) , \(\mathbf{B} = k[\mathbf{X}^4, \mathbf{X}^3\mathbf{Y}, \mathbf{XY}^3, \mathbf{Y}^4]\) and \(\tilde{\mathbf{B}} = k[\mathbf{X}^4, \mathbf{X}^3\mathbf{Y}, \mathbf{X}^2\mathbf{Y}^2, \mathbf{XY}^3, \mathbf{Y}^4]\) of \(\mathbf{C}\) . The inclusions \(\mathbf{A} \subset \mathbf{B} \subset \tilde{\mathbf{B}}\) make it possible to define on \(\mathbf{B}\) and \(\tilde{\mathbf{B}}\) A-module structures. Show that \(\tilde{\mathbf{B}}\) is a flat A-module, while the A-module \(\mathbf{B}\) is not flat.
\item
  Show that an A-module E is flat if and only if it satisfies the following condition (cf. X, p. 15, remark):
\end{enumerate}

(ii'') For every finite family \((c_{i})_{i\in I}\) of elements of A, every solution \(e = (e_i)_{i\in I}\) of the equation \(\sum_{i\in I}c_i e_i = 0\) can be written \(b_1 z_1 + \cdots + b_n z_n\) , where \(b_1, \ldots, b_n \in E\) and where, for \(r = 1, \ldots, n\) , \(z_r = (z_{r,i})_{i \in I}\) is a solution of the equation \(\sum_{i=1}^{n} z_{r,i} c_i = 0\) .

\begin{enumerate}
\def\labelenumi{\arabic{enumi})}
\setcounter{enumi}{6}
\tightlist
\item
  Let E be a monoid and S a subsemigroup of E that is cancellative and satisfies the hypotheses of Exercise 17 of III, p. 121, so that there exists a “monoid of fractions” \(\overline{\mathbf{E}}\) and an injective homomorphism \(\varepsilon: \mathbf{E} \to \overline{\mathbf{E}}\) . Let A be a commutative ring. From \(\varepsilon\) one deduces a homomorphism of monoid algebras \(u: \mathbf{A}^{(\mathbf{E})} \to \mathbf{A}^{(\overline{\mathbf{E}})}\) which allows one to regard \(\mathbf{A}^{\overline{\mathbf{E}}}\) as a \(\mathbf{A}^{(\mathbf{E})}\) -module on the left. Show that this module is flat over \(\mathbf{A}^{(\mathbf{E})}\) (write \(\mathbf{A}^{\overline{\mathbf{E}}}\) as a filtered inductive limit of free \(\mathbf{A}^{(\mathbf{E})}\) -modules of rank 1).
\end{enumerate}

* 8) Let A be the ring of continuous real-valued functions on the interval {[}0, 1{]}.

\begin{enumerate}
\def\labelenumi{\alph{enumi})}
\item
  Show that the ideal of functions vanishing at 0 is a flat A-module (show that it is a filtered inductive limit of the ideals Af, where f is a function which vanishes only at 0).
\item
  Show that the ideal I of functions vanishing in a neighborhood of 0 is a projective A-module (show that there exists a sequence \((f_n)_{n\geqslant 1}\) , \(f_n\in \mathbf{I}\) , such that \(\operatorname {Supp}(f_n)\subset \left[\frac{1}{n + 2},\frac{1}{n}\right]\) and \(\sum_{n}f_{n}(x) = 1\) for \(x\neq 0\) ; deduce
\end{enumerate}

that every element \(h\) of I can be written \(h = \sum_{n}^{t_{h}^{\prime}} (f_n h) e_n\) with \(e_n = \sum_{r \leqslant n+1}^{n} f_r\) , and use II, p. 46, prop. 12 . *

\begin{enumerate}
\def\labelenumi{\arabic{enumi})}
\setcounter{enumi}{8}
\tightlist
\item
  Let P be an A-module. Show that the following conditions are equivalent:
\end{enumerate}

α) For a sequence \(M' \xrightarrow{u} M \xrightarrow{v} M''\) of right A-modules to be exact, it is necessary and sufficient that the sequence

\[
\mathbf {M} ^ {\prime} \otimes_ {\mathbf {A}} \mathbf {P} \xrightarrow {u \otimes 1} \mathbf {M} \otimes_ {\mathbf {A}} \mathbf {P} \xrightarrow {v \otimes 1} \mathbf {M} ^ {\prime \prime} \otimes_ {\mathbf {A}} \mathbf {P}
\]

be exact.

\(\beta)\) P is flat, and for every nonzero right A-module M, one has M \(\otimes_{A} P \neq (0)\) .

γ) P is flat, and for all right A-modules M, N, the canonical homomorphism

\[
\operatorname{Hom} _ {\mathbf {A}} (\mathbf {M}, \mathbf {N}) \rightarrow \operatorname{Hom} _ {\mathbf {Z}} (\mathbf {M} \otimes_ {\mathbf {A}} \mathbf {P}, \mathbf {N} \otimes_ {\mathbf {A}} \mathbf {P})
\]

is injective.

δ) P is flat, and for every maximal right ideal m of A, one has mP ≠ P.

One then says that P is a faithfully flat A-module.

\begin{enumerate}
\def\labelenumi{\arabic{enumi})}
\setcounter{enumi}{9}
\tightlist
\item
  \begin{enumerate}
  \def\labelenumii{\alph{enumii})}
  \tightlist
  \item
    Show that a faithfully flat module is faithful (II, p. 28).
  \end{enumerate}
\end{enumerate}

\begin{enumerate}
\def\labelenumi{\alph{enumi})}
\setcounter{enumi}{1}
\item
  The Z-module Q is faithful and flat, but not faithfully flat.
\item
  Let \((p_n)\) the strictly increasing sequence of prime numbers, and let A be the product ring \(\prod_{n} \mathbf{Z}/p_n \mathbf{Z}\) .
\end{enumerate}

Show that the direct sum E of the \(Z/p_{n}Z\) is a faithful projective A-module which is not faithfully flat (observe that E is an ideal of A such that \(E^{2}=E\) ).

\begin{enumerate}
\def\labelenumi{\arabic{enumi})}
\setcounter{enumi}{10}
\item
  Let \(0 \to \mathbf{P}' \to \mathbf{P} \to \mathbf{P}'' \to 0\) an exact sequence of A-modules. Suppose that \(\mathbf{P}'\) and \(\mathbf{P}''\) are flat, and that one of them is faithfully flat. Show that \(\mathbf{P}\) is faithfully flat.
\item
  Let \(\rho: A \to B\) a ring homomorphism that makes \(B\) an \(A\) faithfully flat right module; let \(M\) be an \(A\) module.
\end{enumerate}

\begin{enumerate}
\def\labelenumi{\alph{enumi})}
\item
  Show that for M to be finitely generated (resp. of finite presentation), it is necessary and sufficient that the same be true of the B-module \(\mathbf{B} \otimes_{\mathbf{A}} \mathbf{M}\) .
\item
  We assume that \(\rho(A)\) is contained in the center of B. Show that M is flat (resp. faithfully flat, resp. finitely generated projective) if and only if the B-module \(B \otimes_{A} M\) is.
\end{enumerate}

\begin{enumerate}
\def\labelenumi{\arabic{enumi})}
\setcounter{enumi}{12}
\tightlist
\item
  Show that every finitely generated flat module P over an integral domain A is projective (consider a presentation \(0 \rightarrow R \xrightarrow{L} L \xrightarrow{P} 0\) , where L is finitely generated free; choose a finitely generated submodule \(R'\) of R such that \(R/R'\) be torsion, and apply Theorem 1, p. 14).
\end{enumerate}

¶ 14) a) Let \(0 \to \mathbf{R} \to \mathbf{L} \to \mathbf{E} \to 0\) an exact sequence of A-modules, where L is a free A-module; let \((e_{a})\) a basis of L. Show that the following conditions are equivalent:

β) For every \(x \in \mathbb{R}\) , if \(\mathfrak{a}_{x}\) is the right ideal generated by the components of \(x\) on the basis \((e_{\alpha})\) , we have \(x \in \mathfrak{a}_{x} \mathbb{R}\) .

\(\gamma)\) For every \(x \in \mathbb{R}\) , there exists a homomorphism \(u_x: \mathbb{L} \to \mathbb{R}\) such that \(u_x(x) = x\) .

δ) For every finite sequence \((x_{i})_{1\leqslant i\leqslant n}\) of elements of R, there exists a homomorphism \(u:L\to R\) such that \(u(x_{i})=x_{i}\) for \(1\leqslant i\leqslant n\) (Use Thm. 1, p. 14.)

\begin{enumerate}
\def\labelenumi{\alph{enumi})}
\setcounter{enumi}{1}
\item
  Let r be the radical of A, and let \(0 \rightarrow R \rightarrow L \rightarrow E \rightarrow 0\) an exact sequence of A-modules, such that L is free. Suppose that E is flat and that R is contained in rL. Show that then R = 0 (with the notations of a), observe that \(a_{x}\) is a finitely generated ideal and that we have \(a_{x} = a_{x} r\) ).
\item
  Let E be a flat finitely generated A-module; suppose there exists a two-sided ideal b of A contained in the radical of A, such that E/bE is a free (A/b)-module. Show that E is then a free A-module (observe that there exists a finitely generated free A-module L such that L/bL is isomorphic to E/bE, and apply b)). In particular, every finitely generated flat module over a local ring is free.
\end{enumerate}

\begin{enumerate}
\def\labelenumi{\arabic{enumi})}
\setcounter{enumi}{14}
\tightlist
\item
  Let A be a ring, \(a\) an element of A. Show that the following properties are equivalent: \(\alpha)\) \(a\in aAa\) .
\end{enumerate}

β) Aa is a direct summand in the module \(A_{s}\) .

γ) \(A_{s}/Aa\) is a flat A-module.

δ) For every right ideal b of A, we have b ∩ Aa = ba.

\begin{enumerate}
\def\labelenumi{\arabic{enumi})}
\setcounter{enumi}{15}
\tightlist
\item
  Let A be a ring. Show that the following properties are equivalent:
\end{enumerate}

α) Every element \(a \in A\) satisfies the equivalent properties of Exercise 15.

β) Every finitely generated left ideal is a direct summand of \(A_{s}\) .

γ) Every left A-module is flat.

δ) Every right A-module is flat.

We then say that A is an absolutely flat ring (cf. VIII, § 8, Exercise 15).

\begin{enumerate}
\def\labelenumi{\arabic{enumi})}
\setcounter{enumi}{16}
\tightlist
\item
  Let A be an absolutely flat ring (Exercise 16).
\end{enumerate}

\begin{enumerate}
\def\labelenumi{\alph{enumi})}
\item
  Let P be a projective A-module. Show that every finitely generated submodule of P is a direct summand of P (reduce to the case where P is finitely generated free).
\item
  Show that every projective A-module is a direct sum of cyclic submodules, isomorphic to cyclic ideals of A. (Use Kaplansky's theorem (II, p. 183, exercise 2) to reduce to the case where P admits a countable family of generators, then use a).)
\item
  Show that every ideal of A generated by a countable family of generators is projective.
\item
  Give an example of an absolutely flat ring A and a finitely generated A-module that is not projective.
\end{enumerate}

18) Let M be an A-module. Show that the following conditions are equivalent:

α) M is finitely generated (resp. finitely presented).

β) For every inductive system \((\mathbf{N}_{i}, u_{ji})\) of A-modules relative to a filtered preordered set I, the canonical homomorphism \(\varinjlim\operatorname{Hom}_{\mathbf{A}}(\mathbf{M}, \mathbf{N}_{i}) \to \operatorname{Hom}_{\mathbf{A}}(\mathbf{M}, \varinjlim \mathbf{N}_{i})\) is injective (resp. bijective).

γ) For every family \((\mathbf{P}_{j})_{j\in J}\) of right A-modules, the canonical homomorphism

\[
\left(\prod_ {j \in J} P _ {j}\right) \otimes_ {A} M \rightarrow \prod_ {j \in J} \left(P _ {j} \otimes_ {A} M\right)
\]

(II, p. 61) is surjective (resp. bijective).

δ) For every set I, the canonical homomorphism \(A_{d}^{1} \otimes_{A} M \to M^{1}\) is surjective (resp. bijective). 19) Let A and B be two rings, \(\varphi : A \to B\) a homomorphism, M a B-module. We equip \(B \otimes_{A} M\) and \(\operatorname{Hom}_{A}(B, M)\) with left B-module structures induced from the B-module structure of B. One says that M is projective relative to A, or \(\varphi\) -projective, if the canonical surjection \(B \otimes_{A} M \to M\) admits a B-linear section; one says that M is injective relative to A, or \(\varphi\) -injective, if the canonical injection

\[
\mathbf {M} \rightarrow \operatorname{Hom} _ {\mathbf {A}} (\mathbf {B}, \mathbf {M})
\]

admits a B-linear retraction.

\begin{enumerate}
\def\labelenumi{\alph{enumi})}
\item
  If M is projective (resp. injective) as an A-module and \(\varphi\) -projective (resp. \(\varphi\) -injective), then M is a projective (resp. injective) B-module.
\item
  For every A-module N, the B-module \(B \otimes_{A} N\) (resp. \(\operatorname{Hom}_{A}(B, N)\) ) is \(\varphi\) -projective (resp. \(\varphi\) -injective).
\item
  We assume that \(\varphi(A)\) is contained in the center of B. Let N be an A-module, P a B-module \(\varphi\) -projective. Show that the B-module \(P \otimes_{A} N\) (resp. the right B-module \(\operatorname{Hom}_{A}(P, N)\) ) is \(\varphi\) -projective (resp. \(\varphi\) -injective).
\end{enumerate}

\begin{enumerate}
\def\labelenumi{\arabic{enumi})}
\setcounter{enumi}{19}
\item
  Let A be a principal ideal ring, \(a\) a nonzero element of A, B = A/Aa. Show that the B-module B \(_{s}\) is injective.
\item
  Show that the following conditions are equivalent:
\end{enumerate}

α) A is Noetherian.

β) Every A-module that is an inductive limit of injective A-modules is injective.

γ) Every A-module that is a direct sum of injective A-modules is injective.

(To prove that \(\gamma\) ) implies \(\alpha\) ), let \((\mathfrak{a}_n)_{n\geqslant 1}\) an increasing sequence of left ideals of A, \(\mathfrak{a} = \bigcup_{n}\mathfrak{a}_{n}\) ,

\(I_{n}\) an injective envelope of \(\mathfrak{a}/\mathfrak{a}_{n}\) . The natural homomorphism from \(\mathfrak{a}\) to \(\oplus I_{n}\) is the restriction of a homomorphism \(f: A \to \oplus I_{n}\) ; consider \(f(1)\) .

\begin{enumerate}
\def\labelenumi{\arabic{enumi})}
\setcounter{enumi}{21}
\tightlist
\item
  Let A be a principal ideal ring, K its field of fractions, \(\pi\) an extremal element of A. Show that the \(\pi\) -primary component of the A-module K/A is an injective envelope of the A-module A/ \(\pi^n\) A for every \(n \geqslant 1\) . ¶ 23) Let I be an injective A-module, E the ring End \(_A\) (I), r the set of elements \(u\) of E such that I is an injective envelope of Ker( \(u\) ).
\end{enumerate}

\begin{enumerate}
\def\labelenumi{\alph{enumi})}
\item
  Show that r is a two-sided ideal and that the ring E/r is absolutely flat (exercise 16). (Let u be an element of E, N a submodule of I such that N ∩ Ker u = 0 and maximal for this property; show that there exists v ∈ E such that vu(x) = x for all x ∈ N, and deduce that uvu - u ∈ r.)
\item
  Show that r is the radical of E (the inclusion r(E) ⊂ r follows from a); show that for every u ∈ r, 1 - u is bijective).
\item
  Assume that I is a direct sum of a finite family of indecomposable injectives. Show that the ring E/r is semisimple.
\end{enumerate}

Let \(i: \mathbf{A} \to \mathbf{I}\) an injective envelope of the A-module \(\mathbf{A}_{s}\) . We set \(\mathbf{E} = \operatorname{End}_{\mathbf{A}}(\mathbf{I})\) and \(\mathbf{B} = \operatorname{End}_{\mathbf{E}}(\mathbf{I})\) , so that \(\mathbf{B}\) is the bicommutant of the A-module I (VIII, § 3, def. 1). We consider \(\mathbf{A}\) as a subring of \(\mathbf{B}\) . One denotes \(j: \mathbf{B} \to \mathbf{I}\) and \(p: \mathbf{E} \to \mathbf{I}\) the maps defined by \(j(b) = b(i(1))\) and \(p(u) = u(i(1))\) for \(b \in \mathbf{B}\) , \(u \in \mathbf{E}\) .

\begin{enumerate}
\def\labelenumi{\alph{enumi})}
\item
  Show that \(p\) is a surjection and \(j\) an injection.
\item
  Show that the following conditions are equivalent:
\end{enumerate}

α) The map p is bijective.

β) The map j is bijective.

γ) The A-module B is injective.

If these conditions are satisfied, the mapping \(j^{-1} \circ p\) is an isomorphism from the ring E onto the opposite ring of B.

\begin{enumerate}
\def\labelenumi{\alph{enumi})}
\setcounter{enumi}{2}
\item
  Let \(\mathfrak{F}\) the set of left ideals \(a\) of \(A\) such that \(a \cap b \neq (0)\) for every nonzero left ideal \(b\) . One denotes \(s(A)\) the set of elements \(a\) of \(A\) for which there exists \(a \in \mathfrak{F}\) such that \(aa = 0\) . Show that \(s(A)\) is a two-sided ideal of \(A\) .
\item
  Show that the following conditions are equivalent:
\end{enumerate}

α) \(s(\mathbf{A}) = 0.\)

β) The ring E is semiprimitive.

γ) The ring B is absolutely flat.

If these conditions are satisfied, then the equivalent properties of \(b\) ) are satisfied.

(Let \(s(I)\) the set of elements \(x\) of I for which there exists \(\mathfrak{a} \in \mathfrak{F}\) such that \(\mathfrak{a} x = 0\) . Using Exercise 23, prove that \(p^{-1}(s(I))\) is equal to the radical \(r\) of E, whence the implication \(\beta) \Rightarrow \alpha\) ). Then show that \(\alpha\) implies \(r = \text{Ker } (p) = 0\) , which entails \(\beta\) and \(\gamma\) (in view of Exercise 23) as well as condition \(\alpha\) of \(b\) ). Finally show that \(\gamma\) implies \(\alpha\) .)

\begin{enumerate}
\def\labelenumi{\alph{enumi})}
\setcounter{enumi}{4}
\item
  If A is left Noetherian, the ideal s(A) is nilpotent (first prove that every element a of s(A) is nilpotent, by considering the left annihilators of the \(a^{n}\) ; then use Exercise 8, b) of VIII, § 9). ¶ 25) Assume that the ring A is left Noetherian and contains no nonzero nilpotent two-sided ideals. a) Prove that the ring B defined in Exercise 24 is semisimple (“Goldie's theorem”; use Exercise 23, c) and Exercise 24, d) and e)). If moreover (0) is a prime two-sided ideal in A (VIII, § 8, Exercise 19), show that B is simple.
\item
  One identifies A with a subring of B. Prove the following properties:
\end{enumerate}

\begin{enumerate}
\def\labelenumi{(\roman{enumi})}
\item
  Every element of A that is left cancellable is invertible in B (hence cancellable in A).
\item
  Every element \(b\) of \(\mathbf{B}\) can be written \(s^{-1}a\) , with \(s \in \mathbf{A}\) , \(a \in \mathbf{A}\) and \(s\) cancellable.
\end{enumerate}

One sometimes says that B is the total ring of left fractions of A.

(To prove (ii), reduce to the case where B is simple. Show then that the ideal (0) of A is prime. Let a be an ideal of F, a an element of a for which the right ideal aB is maximal among the ideals a'B (a' ∈ a), and x an element of B such that axa = a. Prove that B(1 - xa) ⊂ aB, whence

\[
(\mathbf {B} (1 - a x) \cap \mathfrak {a}) \cdot (\mathbf {B} (1 - x a) \cap \mathfrak {a}) = (0);
\]

deduce that \(a\) is cancellable. Conclude by taking for a the set of all \(a \in A\) such that \(ab \in A\) .) c) Assume that every nonzero element of \(A\) is cancellable. Show that \(B\) is a field, which is a left field of fractions for \(A\) (cf. I, p. 155, Exercise 15).

¶ 26) a) Show that the following conditions are equivalent:

α) The ring A is right Noetherian, and the A-module \(A_{s}\) is injective.

\(\alpha^{0}\) A is left Noetherian and the A-module \(A_{d}\) is injective.

β) The ring A is self-injective (VIII, § 2, Exercise 11).

(Prove that \(\alpha\) ) implies that every right ideal r satisfies \(r = r^0\) . Then show that \(r(A)\) is nilpotent, and deduce that A is right Artinian using Exercise 23. Next prove that the dual of a simple left A-module is simple. Deduce that A is left Noetherian and finally self-injective. b) Show that the ring A is self-injective if and only if every injective A-module is projective and every projective A-module is injective.

¶ 27) Assume the ring A is left Noetherian.

\begin{enumerate}
\def\labelenumi{\alph{enumi})}
\item
  Let I be an indecomposable injective A-module. Show that the set of annihilator ideals of a nonzero submodule of I admits a greatest element \(\mathfrak{A}(I)\) , which is a prime two-sided ideal (VIII, § 8, Exercise 19).
\item
  Show that for every prime two-sided ideal p of A, there exists an indecomposable injective A-module I such that \(\mathfrak{A}(I) = \mathfrak{p}\) (consider an indecomposable direct factor of the injective envelope of A/p).
\item
  Assume A is commutative. Show that every indecomposable injective A-module I is the injective envelope of A/U(I). Deduce that the map I ↦ U(I) defines a bijection from the set S of isomorphism classes of indecomposable injective A-modules onto the set of prime ideals of A. The inverse bijection associates to the ideal p the injective envelope of A/p.
\item
  Show that for the map I ↦ \(\mathfrak{A}(I)\) from \(\mathfrak{A}\) onto the set of prime two-sided ideals of A to be bijective, it suffices that A satisfy the following condition (H):
\end{enumerate}

\begin{enumerate}
\def\labelenumi{(\Alph{enumi})}
\setcounter{enumi}{7}
\tightlist
\item
  For every left ideal a of A, there exist elements \(x_{1}, \ldots, x_{k}\) of A/a such that
\end{enumerate}

\[
\operatorname{Ann} (\mathbf {A} / \mathfrak {a}) = \bigcap_ {i} \operatorname{Ann} (x _ {i}).
\]

\begin{enumerate}
\def\labelenumi{\alph{enumi})}
\setcounter{enumi}{4}
\tightlist
\item
  Show that condition (H) is satisfied in the following cases:
\end{enumerate}

α) All left ideals of A are two-sided.

β) A is Artinian.

γ) The center Z of A is a Noetherian ring and A is a finitely generated Z-module.

\begin{enumerate}
\def\labelenumi{\alph{enumi})}
\setcounter{enumi}{5}
\tightlist
\item
  Let k be a field, \(Q^{+}\) the additive monoid of rational numbers \(\geqslant 0\) , A the algebra \(k^{(\mathbf{Q}^{+})}\) , a a nonzero element of A, I the injective envelope of A/Aa. Show that the A-module I is indecomposable, but contains no submodule isomorphic to A/p for p prime (note that the set of ideals of A is totally ordered and that A contains only two prime ideals).
\end{enumerate}

\begin{enumerate}
\def\labelenumi{\arabic{enumi})}
\setcounter{enumi}{27}
\tightlist
\item
  Let A be a Noetherian ring, M an A-module, I an injective envelope of M, a direct sum of a family \((I_{\alpha})_{\alpha \in J}\) of indecomposable submodules. We denote by Ass (M) the set of prime two-sided ideals \(\mathfrak{U}(I_{\alpha})\) , for \(\alpha \in J\) (exercise 27); if \(p \in \operatorname{Ass}(M)\) , we say that p is associated with M.
\end{enumerate}

\begin{enumerate}
\def\labelenumi{\alph{enumi})}
\item
  We have Ass (M) ≠ ∅ if M ≠ 0; if M is finitely generated, the set Ass (M) is finite.
\item
  Show that a prime two-sided ideal p is associated with M if and only if M contains a submodule N such that Ann (N') = p for every nonzero submodule N' of N. If A is commutative, p is associated with M if and only if M contains a submodule isomorphic to A/p.
\item
  Let M' be a submodule of M. Show that we have
\end{enumerate}

\[
\operatorname{Ass} \left(\mathbf {M} ^ {\prime}\right) \subset \operatorname{Ass} (\mathbf {M}) \subset \operatorname{Ass} \left(\mathbf {M} ^ {\prime}\right) \cup \operatorname{Ass} \left(\mathbf {M} / \mathbf {M} ^ {\prime}\right).
\]

\begin{enumerate}
\def\labelenumi{\alph{enumi})}
\setcounter{enumi}{3}
\tightlist
\item
  We set \(J_{p} = \sum_{\mathfrak{M}(I_{\alpha}) = p} I_{\alpha}\) and \(Q(p) = \left( \sum_{q \neq p} J_{q} \right) \cap M\) Show that we have \(\bigcap_{p \in \operatorname{Ass}(M)} Q(p) = 0\) and \(\bigcap_{p \in E} Q(p) \neq 0\) for every subset E of Ass(M) distinct from Ass(M). For every \(p \in \operatorname{Ass}(M)\) , we have
\end{enumerate}

Ass \((\mathbf{M} / \mathbf{Q}(\mathfrak{p})) = \{\mathfrak{p}\}\) .

\begin{enumerate}
\def\labelenumi{\alph{enumi})}
\setcounter{enumi}{4}
\tightlist
\item
  Let a be an element of the center of A. Show that for the homothety of ratio a to be injective in M, it is necessary and sufficient that a belong to none of the prime ideals associated with M.
\end{enumerate}

Assume that the ring A is commutative and Noetherian. Let a be an ideal of A.

\begin{enumerate}
\def\labelenumi{\alph{enumi})}
\tightlist
\item
  Let H be an A-module such that every element of H is annihilated by a power of a. Show that for H to be injective, it suffices that for every integer \(n \geqslant 1\) every finitely generated A-module M annihilated by \(a^{n}\) and every submodule \(M'\) of M, the natural homomorphism \(\operatorname{Hom}_{\mathbf{A}}(\mathbf{M}, \mathbf{H}) \to \operatorname{Hom}_{\mathbf{A}}(\mathbf{M}', \mathbf{H})\) be surjective.
\end{enumerate}

(Use the fact that for every ideal b of A, there exists an integer \(k \geqslant 0\) such that \(\mathfrak{b} \cap \mathfrak{a}^k \subset \mathfrak{ba}\) (AC, III, § 3, no. 1, cor. 2).)

\begin{enumerate}
\def\labelenumi{\alph{enumi})}
\setcounter{enumi}{1}
\item
  Let M be an injective A-module. For \(n \geqslant 1\) , we denote by \(\mathbf{H}_{n}\) the submodule of M consisting of the elements annihilated by \(\alpha^{n}\) ; set \(\mathbf{H} = \bigcup \mathbf{H}_{n}\) . Show that H is an injective A-module.
\item
  Show that every element of the injective envelope of A/a is annihilated by a power of a. * 30) Assume the ring A is local, commutative, and Noetherian; denote its maximal ideal by m, and k = A/m. Let I be an A-module such that every element of I is annihilated by a power of m. For every A-module M, denote by T(M) the A-module Hom\_A (M, I) and c\_M : M → T(T(M)) the homomorphism defined by
\end{enumerate}

\[
\left(c _ {\mathbf {M}} (m)\right) (u) = u (m) \quad \text { for } \quad m \in \mathbf {M}  , \quad u \in \mathrm{T} (\mathbf {M})  .
\]

Show that the following properties are equivalent:

α) For every finite-length A-module M, the A-module T(M) is of finite length and the homomorphism \(c_{M}: M \to T(T(M))\) is bijective.

β) For every A-module M of finite length, one has long T(M) = long M.

γ) The A-module I is injective, and the A-modules k and T(k) are isomorphic.

δ) The A-module I is isomorphic to the injective envelope of k.

When these conditions are satisfied, one sometimes says that I is a dualizing module for the ring A. (Show that property γ) is equivalent to each of the other three: to see that α) or β) implies γ), use Exercise 29.)

\begin{enumerate}
\def\labelenumi{\arabic{enumi})}
\setcounter{enumi}{30}
\item
  With the hypotheses of the previous exercise, let B be a local, commutative, Noetherian ring, and let \(\rho: A \to B\) a ring homomorphism that makes B a finitely generated A-module. If I is a dualizing module for A, show that the B-module \(\mathrm{Hom}_{\mathbf{A}}(\mathbf{B}, I)\) is dualizing for B. In particular, for every ideal \(a\) of A, the submodule of I consisting of elements annihilated by \(a\) is a dualizing module for A/a.
\item
  \begin{enumerate}
  \def\labelenumii{\alph{enumii})}
  \tightlist
  \item
    We keep the hypotheses of Exercise 30; we further assume that A contains a field \(k_0\) such that \(k\) is a finite-dimensional \(k_0\) -vector space. Show that the A-module \(l = \varinjlim_{n} \operatorname{Hom}_{k_0}(A/m^n, k_0)\)
  \end{enumerate}
\end{enumerate}

is dualizing for A.

\begin{enumerate}
\def\labelenumi{\alph{enumi})}
\setcounter{enumi}{1}
\tightlist
\item
  Let \(k\) a field of characteristic zero, B the ring of formal power series \(k[[\mathrm{T}_{1}, \ldots, \mathrm{T}_{n}]]\) . We endow the group \(\mathbf{I} = k[\mathrm{X}_{1}, \ldots, \mathrm{X}_{n}]\) with a B-module structure by setting \(\mathrm{T}_{i} \cdot \mathbf{P} = \partial \mathbf{P} / \partial \mathbf{X}_{i}\) for all \(\mathbf{P} \in k[\mathrm{X}_{1}, \ldots, \mathrm{X}_{n}]\) . Show that I is an injective envelope of \(k\) .
\end{enumerate}

\exercisesection{Complexes of A-Modules}

\begin{enumerate}
\def\labelenumi{\arabic{enumi})}
\tightlist
\item
  Assume the ring A is principal; let C be a complex of free A-modules.
\end{enumerate}

\begin{enumerate}
\def\labelenumi{\alph{enumi})}
\item
  Show that C is the direct sum of a family of complexes \((^{p}\mathbf{E})_{p\in \mathbb{Z}}\) satisfying: \(^{p}\mathbf{E}_{n} = 0\) for \(n\neq p\) , \(p - 1 ; Z_{n}(^{p}\mathbf{E}) = 0\) .
\item
  We further assume that \(C_n\) is finitely generated for all \(n \in \mathbf{Z}\) . Show that \(C\) is a direct sum of complexes having only either one nonzero component, or two consecutive nonzero components, these components being isomorphic to \(A\) .
\end{enumerate}

\begin{enumerate}
\def\labelenumi{\arabic{enumi})}
\setcounter{enumi}{1}
\item
  Let C and C' be two complexes of A-modules; suppose that the complexes C and B(C) (resp. C' and Z(C')) are projective (resp. injective). Show that every morphism of complexes from H(C) to H(C') is induced by a morphism from C to C'.
\item
  Let C be a complex of A-modules; one regards C as an A(ε)-module (X, p. 27). We denote by
\end{enumerate}

\[
i: \mathbf {A} \rightarrow \mathbf {A} (\varepsilon)
\]

the canonical injection.

Show that the following conditions are equivalent:

\begin{enumerate}
\def\labelenumi{(\roman{enumi})}
\item
  C is i-projective (X, p. 170, exercise 19).
\item
  It is i-injective (loc. cit.).
\item
  The complex C is homotopic to zero.
\item
  There exists a graded A-module M and an isomorphism of graded A(ε)-modules A(ε) ⊗A M → C.
\item
  There exists a graded A-module N and an isomorphism of graded A(ε)-modules C → Hom\_A(A(ε), N).
\end{enumerate}

\begin{enumerate}
\def\labelenumi{\arabic{enumi})}
\setcounter{enumi}{3}
\tightlist
\item
  \begin{enumerate}
  \def\labelenumii{\alph{enumii})}
  \tightlist
  \item
    Let C be a right-bounded A-complex, such that C and H(C) are projective. Show that C is split.
  \end{enumerate}
\end{enumerate}

\begin{enumerate}
\def\labelenumi{\alph{enumi})}
\setcounter{enumi}{1}
\tightlist
\item
  On the ring B = A(ε) (X, p. 27), we consider the complex C defined by C\_n = B for all n, and d(b) = εb for b ∈ B. The complex C is free with zero homology, but not split.
\end{enumerate}

\begin{enumerate}
\def\labelenumi{\arabic{enumi})}
\setcounter{enumi}{4}
\tightlist
\item
  Let C be a complex, and let s be a splitting of C, \(\delta: \mathbf{C} \to \mathbf{C}\) a graded homomorphism of degree -1, such that \((d + \delta) \circ (d + \delta) = 0\) . One denotes \(\mathbf{C}'\) the complex \((\mathbf{C}, d + \delta)\) . Assume that the A-linear map \(\theta\) from C to C defined by \(\theta = 1_C + \delta s\) is bijective.
\end{enumerate}

\begin{enumerate}
\def\labelenumi{\alph{enumi})}
\item
  In order for \(s\theta^{-1}\) to be a splitting of \(C'\) , it is necessary and sufficient that \(\delta(Z(C)) \subset \theta(B(C))\) .
\item
  Assume the conditions of a) are satisfied. Show that there exist direct sum decompositions: \(\mathbf{C} = \mathbf{B}(\mathbf{C}) \oplus \mathbf{Ker}(ds) = \mathbf{B}(\mathbf{C}') \oplus \mathbf{Ker}(ds)\) . If p (resp. \(p'\) ) is the projector with image \(\mathbf{Ker}(ds)\) and kernel \(\mathbf{B}(\mathbf{C})\) (resp. \(\mathbf{B}(\mathbf{C}')\) ), we have \(p' = p\theta^{-1}\) .
\item
  We further assume that the A-linear map \(\lambda\) from C to C defined by \(\lambda = 1_{C} + s\delta\) is bijective. Show that we have \(Z(C') = \lambda^{-1}(Z(C))\) and \(C = Z(C) \oplus \operatorname{Im}(sd) = Z(C') \oplus \operatorname{Im}(sd)\) . If \(q\) (resp. \(q'\) ) is the projector with image \(Z(C)\) (resp. \(Z(C')\) ) and kernel \(\operatorname{Im}(sd)\) , we have \(q' = \lambda^{-1} q\) .
\end{enumerate}

\begin{enumerate}
\def\labelenumi{\arabic{enumi})}
\setcounter{enumi}{5}
\tightlist
\item
  We keep the notation from Exercise 5; we further assume that A is commutative, that the complexes C and \(\mathbf{C}^{\prime}/\mathbf{B}(\mathbf{C}^{\prime})\) are flat and that there exists an ideal m of A satisfying \(\delta(\mathbf{C}) \subset \mathfrak{m}\mathbf{C}\) and \(\bigcap_{n} (\mathbf{B}(\mathbf{C}) + \mathfrak{m}^{n}\mathbf{C}) = \mathbf{B}(\mathbf{C})\) (This last condition is always satisfied if A is a local Noetherian ring with maximal ideal m and if \(C_{n}\) is a finitely generated A-module for every n, cf. AC, III, § 3, n° 3, prop. 6). Then show that \(s\theta^{-1}\) is a splitting of \(C^{\prime}\) .
\end{enumerate}

(The flatness hypotheses imply \(\operatorname{Im}(d + \delta) \cap \mathfrak{m}^n\mathbf{C} = (d + \delta)\, (\mathfrak{m}^n\,\mathbf{C})\) for all \(n\) ; if \(x \in \mathbf{Z}(\mathbf{C})\) , construct by induction elements \(y_n \in \mathbf{C}\) , \(z_n \in \mathbf{Z}(\mathbf{C}) \cap \mathfrak{m}^n\mathbf{C}\) such that \(\delta(x) = \theta d(y_n) + \delta(z_n)\) . Conclude using exercise 5, a).)

\begin{enumerate}
\def\labelenumi{\arabic{enumi})}
\setcounter{enumi}{6}
\tightlist
\item
  Let C and C' be two complexes, \(s_{1}(\text{resp. } s')\) a splitting of C (resp. C'), \(u : C' \to C\) a morphism of complexes. We define an A-endomorphism \(\sigma\) of Con(u), graded of degree -1, by setting
\end{enumerate}

\[
\sigma (y ^ {\prime}, x) = \left(- s ^ {\prime} (y ^ {\prime}), s (x) - s u s ^ {\prime} (y ^ {\prime})\right).
\]

Show that in order for \(\sigma\) to be a splitting of Con \((u)\) , it is necessary and sufficient that we have H(u) = 0.

\begin{enumerate}
\def\labelenumi{\arabic{enumi})}
\setcounter{enumi}{7}
\tightlist
\item
  Let \(0 \to M' \xrightarrow{\mu} M \xrightarrow{\nu} M'' \to 0\) an exact sequence of A-modules, which can be viewed as a sequence
\end{enumerate}

exact sequence of complexes (zero in degrees ≠ 0). Show that the morphism of complexes φ : Con(u) → M''\,(X, p. 39) is a homotopy equivalence if and only if the exact sequence 0 → M' → M → M'' → 0 is split. 9) Let u : C' → C be a morphism of A-complexes. Denote by v : C → Coker(u) the quotient map. Consider the graded A-homomorphisms

\[
\alpha : (\operatorname{Ker} u) (- 1) \rightarrow \operatorname{Con} (u) \quad \text { and } \quad \varphi : \operatorname{Con} (u) \rightarrow \operatorname{Coker} (u)
\]

defined by \(\alpha(x') = (x', 0)\) , \(\varphi(x', x) = v(x)\) for \(x' \in C'\) , \(x \bullet C\) . Show that \(\alpha, \varphi\) are morphisms of complexes and that they give rise to an exact sequence

\[
\dots \rightarrow \mathrm{H} _ {n - 1} (\text { Ker } u) \xrightarrow {\mathrm{H} (\alpha)} \mathrm{H} _ {n} (\text { Con } (u)) \xrightarrow {\mathrm{H} (\varphi)} \mathrm{H} _ {n} (\text { Coker } u) \xrightarrow {\delta} \mathrm{H} _ {n - 2} (\text { Ker } u) \rightarrow \dots
\]

where δ is the composite of the connecting homomorphisms relative to the two exact sequences

\[
0 \rightarrow \operatorname{Ker} u \rightarrow C ^ {\prime} \rightarrow \operatorname{Im} u \rightarrow 0 \quad \text { and } \quad 0 \rightarrow \operatorname{Im} u \rightarrow C \rightarrow \operatorname{Coker} u \rightarrow 0.
\]

\begin{enumerate}
\def\labelenumi{\arabic{enumi})}
\setcounter{enumi}{9}
\tightlist
\item
  Let \(u: \mathbf{C} \to \mathbf{C}'\) a morphism of A-complexes, \(\mathbf{P}\) a projective A-module, \(n\) be an integer \(\geqslant 0\) ,
\end{enumerate}

\[
h: \mathrm{P} \rightarrow \mathrm{H} _ {n} (\text { Con } (u))
\]

an A-homomorphism. Show that there exists a complex \(\tilde{C}\) , an exact sequence of complexes

\[
0 \to \mathbf {C} \xrightarrow {i} \tilde {\mathbf {C}} \to \mathbf {P} (- n) \to 0
\]

and a morphism \(\tilde{u}:\tilde{\mathbf{C}}\to \mathbf{C}'\) such that \(\tilde{u}\circ i = u\) and that the morphism \(\mathbf{C}(i):\mathrm{Con}(u)\to \mathrm{Con}(\tilde{u})\) induced by \(i\) satisfies the following properties:

α) The map \(\mathrm{H}_{p}(\mathrm{C}(i)) : \mathrm{H}_{p}(\mathrm{Con}(u)) \to \mathrm{H}_{p}(\mathrm{Con}(\tilde{u}))\) is bijective for \(p \neq n, n + 1\) ;

β) There exists an exact sequence

\[
0 \rightarrow \mathrm{H} _ {n + 1} (\operatorname{Con} (u)) \xrightarrow {\mathrm{H} _ {n + 1} (\mathbf {C} (i))} \mathrm{H} _ {n + 1} (\operatorname{Con} (\tilde {u})) \rightarrow \mathrm{P} \xrightarrow {h} \mathrm{H} _ {n} (\operatorname{Con} (u)) \xrightarrow {\mathrm{H} _ {n} (\mathbf {C} (i))} \mathrm{H} _ {n} (\operatorname{Con} (\tilde {u})) \rightarrow 0.
\]

\begin{enumerate}
\def\labelenumi{\arabic{enumi})}
\setcounter{enumi}{10}
\tightlist
\item
  We call a bicomplex of A-modules a triple (C, \(d'\) , \(d''\) ), where C is a graded A-module of type \(\mathbf{Z} \times \mathbf{Z}\) , \(d': \mathbf{C} \to \mathbf{C}\) a graded A-endomorphism of degree (-1, 0) and \(d'': \mathbf{C} \to \mathbf{C}\) a graded A-endomorphism of degree (0, -1), satisfying \(d' \circ d' = d'' \circ d'' = d' \circ d'' + d'' \circ d' = 0\) We will also use the ascending graduation defined by \(\mathbf{C}^{p,q} = \mathbf{C}_{-p,-q}\) for \(p, q \in \mathbf{Z}\) .
\end{enumerate}

We denote by \(Z'(C)\) , \(Z''(C)\) , \(B'(C)\) , \(B''(C)\) the sub-bicomplexes \(\text{Ker } (d')\) , \(\text{Ker } (d'')\) , \(\text{Im } (d')\) , \(\text{Im } (d'')\) , and \(H'(C)\) (resp. \(H''(C)\) ) the bicomplex \(Z'(C)/B'(C)\) (resp. \(Z''(C)/B''(C)\) ).

\begin{enumerate}
\def\labelenumi{\alph{enumi})}
\tightlist
\item
  We set \(d = d' + d''\) . Show that when C is equipped with the total graduation associated with the given bigraduation (defined by \(\mathbf{C}_n = \bigoplus_{p+q=n} \mathbf{C}_{p,q}\) , cf. II, p. 164), the pair (C, d) is a complex. This is called
\end{enumerate}

the complex associated with the bicomplex (C, \(d'\) , \(d''\) ); its homology is denoted H(C).

\begin{enumerate}
\def\labelenumi{\alph{enumi})}
\setcounter{enumi}{1}
\item
  Let (C, \(d'\) , \(d''\) ) and (K, \(\delta'\) , \(\delta''\) ) be two bicomplexes. A morphism from (C, \(d'\) , \(d''\) ) to (K, \(\delta'\) , \(\delta''\) ) is a graded A-homomorphism \(u\) of degree (0, 0) from C to K satisfying: \(\delta'\circ u = u \circ d'\) and \(\delta''\circ u = u \circ d''\) . Show that \(u\) induces a morphism of the associated complexes.
\item
  Let \(u, v\) two morphisms of bicomplexes from C to K. One calls a homotopy connecting \(u\) to \(v\) a pair \((s', s'')\) of A-homomorphisms from C to K, graded of degrees (1, 0) and (0, 1) respectively, such that
\end{enumerate}

\[
g - f = d ^ {\prime} \circ s ^ {\prime} + s ^ {\prime} \circ d ^ {\prime} + d ^ {\prime \prime} \circ s ^ {\prime \prime} + s ^ {\prime \prime} \circ d ^ {\prime \prime}; \quad s ^ {\prime} \circ d ^ {\prime \prime} + d ^ {\prime \prime} \circ s ^ {\prime} = 0; \quad s ^ {\prime \prime} \circ d ^ {\prime} + d ^ {\prime} \circ s ^ {\prime \prime} = 0.
\]

Show that \(s' + s''\) is a homotopy connecting the associated complex morphisms induced by \(u\) and \(v\) . 12) Let \((\mathbf{C}, d', d'')\) a bicomplex, such that one has \(\mathrm{H}_{p,-p}''(\mathbf{C}) = 0\) for all \(p \in \mathbf{Z}\) , \(\mathrm{H}_{-q,q+1}'(\mathbf{C}) = 0\) for \(q \geqslant 0\) , \(\mathrm{H}_{r,-r-1}'(\mathbf{C}) = 0\) for \(r \geqslant 0\) , and that there exist positive integers \(a\) and \(b\) such that \(\mathbf{C}_{a,-a} = \mathbf{C}_{-b,b} = 0\) . Prove that the A-module \(\mathrm{H}_{0,0}'(\mathbf{C})\) is zero.

\begin{enumerate}
\def\labelenumi{\arabic{enumi})}
\setcounter{enumi}{12}
\tightlist
\item
  Let I be a finite set; we denote by \((e_i)_{i\in I}\) the canonical basis of the Z-module \(\mathbf{Z}^{(I)}\) . We call an I-complex (resp. I-precomplex) of A-modules a graded A-module C of type \(\mathbf{Z}^{I}\) , equipped with a family of A-endomorphisms \((d_i)_{i\in I}\) , such that \(d_i\) be graded of degree \((-e_i)\) satisfying:
\end{enumerate}

\[
d _ {i} \circ d _ {i} = 0 \quad \text { for every } \quad i \in I;
\]

\[
d _ {i} \circ d _ {j} + d _ {j} \circ d _ {i} = 0 (\text { resp. } d _ {i} \circ d _ {j} = d _ {j} \circ d _ {i}) \quad \text { for } \quad i, j \in I.
\]

If Card (I) = 1, an I-complex is a complex; if Card (I) = 2, an I-complex is a bicomplex (exercise 11). If \((\mathbf{C}, (d_i))\) and \((\mathbf{C}', (d_i'))\) are two I-complexes (resp. I-precomplexes), a morphism of \((\mathbf{C}, (d_i))\) to \((\mathbf{C}', (d_i'))\) is a graded homomorphism of degree zero \(u: \mathbf{C} \to \mathbf{C}'\) such that \(d_i' \circ u = u \circ d_i\) for all \(i \in I\) .

\begin{enumerate}
\def\labelenumi{\alph{enumi})}
\tightlist
\item
  Let (C, \((d_i)\) ) an I-complex. If one equips C with the total grading (of type Z) deduced from its grading, show that (C, \(\sum_{i\in I}d_i\) ) is a complex of A-modules, called the complex associated with the I-complex (C, \((d_i)\) ).
\end{enumerate}

Show that a morphism of I-complexes induces a morphism of the associated complexes.

\begin{enumerate}
\def\labelenumi{\alph{enumi})}
\setcounter{enumi}{1}
\tightlist
\item
  Let C be an I-precomplex, and let R be a total order relation on I, denoted ≤. For \(k = (k_i)_{i \in I}, x \in C_k\) and \(i \in I\) , we set
\end{enumerate}

\(\delta_{i}(x) = (-1)^{\sum_{j < i} k_{j}} d_{i}(x)\) . Show that \((C, (\delta_{i}))\) is an I-complex, denoted \(C_{R}\) .

\begin{enumerate}
\def\labelenumi{\alph{enumi})}
\setcounter{enumi}{2}
\item
  Let \(\mathbf{R}_{\alpha}, \mathbf{R}_{\beta}\) two total order relations on I, denoted \(\leqslant_{\alpha}\) and \(\leqslant_{\beta}\) . Let \(u_{\beta \alpha}\) the endomorphism of C defined by \(u(x) = (-1)^{\varepsilon(k)} x\) for \(x \in \mathbf{C}_k\) , with \(\varepsilon(k) = \sum_{\substack{i < _{\alpha} j \\ j < _{\beta} i}} k_i k_j\) . Show that \(u_{\beta \alpha}\) is a morphism of I-complexes of \(\mathbf{C}_{\mathbf{R}_{\alpha}}\) to \(\mathbf{C}_{\mathbf{R}_{\beta}}\) .
\item
  Let S be the set of total order relations on I; we endow S with the trivial preorder (that is, the preorder such that \(R_{\alpha} \leqslant R_{\beta}\) whatever \(R_{\alpha}, R_{\beta} \in S\) ). Show that the system \((\mathbf{C}_{R\alpha}, u_{\alpha\beta})\) is a projective system of I-complexes, relative to S; its limit C is an I-complex, called the I-complex associated with the I-precomplex C. Show that every morphism of I-precomplexes induces a morphism of the associated I-complexes.
\end{enumerate}

A spectral sequence of A-modules is defined as a sequence of A-complexes \((\mathrm{E}_r, d_r)_{r \geqslant 2}\) and of mappings \(\alpha_r: \mathrm{H}(\mathrm{E}_r) \to \mathrm{E}_{r+1}\) such that:

\begin{enumerate}
\def\labelenumi{(\roman{enumi})}
\tightlist
\item
  The grading of the complex \(E_r\) is the total grading associated with a bigrading
\end{enumerate}

\[
\mathrm{E} _ {r} = \bigoplus_ {p, q \in \mathbf {Z}} \mathrm{E} _ {r} ^ {p, q} (r \geqslant 2);
\]

we have \(d_r(E_r^{pq}) \subset E_r^{p + r,q - r + 1}\) .

\begin{enumerate}
\def\labelenumi{(\roman{enumi})}
\setcounter{enumi}{1}
\tightlist
\item
  The maps \(\alpha_{r}\) are isomorphisms of bigraded A-modules.
\end{enumerate}

We will also use the descending bigrading of \(E_{r}\) : we denote \(E_{pq}^{r} = E_{r}^{-p, -q}\) .

\begin{enumerate}
\def\labelenumi{\alph{enumi})}
\tightlist
\item
  Here we denote by \(Z_{r+1}(E_r)\) (resp. \(B_{r+1}(E_r)\) ) the bigraded A-module Ker \(d_r\) (resp. Im \(d_r\) ). Let
\end{enumerate}

\[
p _ {r}: \mathrm{E} _ {r} \rightarrow \mathrm{E} _ {r} / \mathrm{B} _ {r + 1} (\mathrm{E} _ {r})
\]

the quotient map and \(i_{r}: \mathbf{E}_{r+1} \to \mathbf{E}_{r}/\mathbf{B}_{r+1}(\mathbf{E}_{r})\) the composite map of \(\alpha_{r}^{-1}\) and the natural injection. For every \(r \geqslant 2\) , we define by induction on \(k\) the bigraded submodules \(\mathbf{Z}_{r+k}(\mathbf{E}_{r})\) and \(\mathbf{B}_{r+k}(\mathbf{E}_{r})\) of \(\mathbf{E}_{r}\) by setting:

\[
\mathbf {Z} _ {r + k} (\mathrm{E} _ {r}) = p _ {r} ^ {- 1} \left(i _ {r} \left(\mathbf {Z} _ {r + k} (\mathrm{E} _ {r + 1})\right)\right); \quad \mathbf {B} _ {r + k} (\mathrm{E} _ {r}) = p _ {r} ^ {- 1} \left(i _ {r} \left(\mathbf {B} _ {r + k} (\mathrm{E} _ {r + 1})\right)\right).
\]

Prove that we have the inclusions:

\[
0 \subset \mathrm{B} _ {r + 1} (\mathrm{E} _ {r}) \subset \mathrm{B} _ {r + 2} (\mathrm{E} _ {r}) \subset \dots \subset \mathrm{Z} _ {r + 2} (\mathrm{E} _ {r}) \subset \mathrm{Z} _ {r + 1} (\mathrm{E} _ {r}) \subset \mathrm{E} _ {r}
\]

and that \(\mathbf{E}_k\) is isomorphic to \(\mathbf{Z}_k(\mathbf{E}_r) / \mathbf{B}_k(\mathbf{E}_r)\) for \(k\geqslant r + 1\)

We set \(Z_{\infty}(E_2) = \bigcap_{r \geqslant 3} Z_r(E_2), B_{\infty}(E_2) = \bigcup_{r \geqslant 3} B_r(E_2)\) and \(E_{\infty} = Z_{\infty}(E_2)/B_{\infty}(E_2)\) . We say that the spectral sequence is regular if the decreasing sequence \((Z_r(E_2^{pq}))_{r \geqslant 2}\) is stationary for every pair \((p, q)\) .

\begin{enumerate}
\def\labelenumi{\alph{enumi})}
\setcounter{enumi}{1}
\item
  Let \(E = (E_r, d_r, \alpha_r)_{r \geqslant 2}\) and \(E' = (E_r', d_r', \alpha_r')_{r \geqslant 2}\) two spectral sequences of A-modules; a morphism from E to \(E'\) is a sequence of morphisms of A-complexes \(u_r : E_r \to E_r'\) , compatible with the bigradings, such that \(\alpha_r' \circ H(u_r) = u_{r+1} \circ \alpha_r\) for all \(r \geqslant 2\) Show that for every \(r \geqslant 2\) and for every \(k \geqslant r + 1\) , \(u_r\) induces a homomorphism of \(Z_k(E_r)\) (resp. \(B_k(E_r)\) in \(Z_k(E_r')\) (resp. \(B_k(E_r')\) ). In particular, \(u_2\) induces homomorphisms \(Z_\infty(u_2) : Z_\infty(E_2) \to Z_\infty(E_2')\) , \(B_\infty(u_2) : B_\infty(E_2) \to B_\infty(E_2')\) and \(u_\infty : E_\infty \to E_\infty'\) ; all these homomorphisms are compatible with the bigradings.
\item
  We assume that \(u_2\) is an isomorphism. Show that then \(u_r\) is an isomorphism for all \(r \geqslant 2\) as well as \(\mathbf{B}_{\infty}(u_2)\) . If moreover the sequences E and E' are regular, \(\mathbf{Z}_{\infty}(u_2)\) and \(u_{\infty}\) are isomorphisms.
\item
  Let \(G = \bigoplus_{n \in Z} G^n\) a graded A-module, equipped with a decreasing sequence \((F^p G)_{p \in \mathbb{Z}}\) of graded sub-A-modules
\end{enumerate}

such that \(\bigcap_{p} F^{p} G = 0\) and \(\bigcup_{p} F^{p} G = G\) . We say that the spectral sequence E approaches \((G, (F^{p} G)_{p \in \mathbb{Z}})\)

(or simply G) if there exist isomorphisms \(\beta^{pq}: E_{\infty}^{pq} \to (F^p G)^{p+q}/(F^{p+1} G)^{p+q}\) for all

\[
(p, q) \in \mathbf {Z} \times \mathbf {Z}.
\]

One says that the spectral sequence E converges to G if it approaches G and if moreover:

\begin{enumerate}
\def\labelenumi{(\roman{enumi})}
\item
  The spectral sequence E is regular;
\item
  For every integer n, there exists an integer p such that \((\mathbf{F}^{p}\mathbf{G})^{n}=0\) .
\end{enumerate}

Suppose that E (resp. E′) converges to G (resp. G′). Let \(u: \mathrm{E} \to \mathrm{E}'\) a morphism of spectral sequences, \(v: \mathrm{G} \to \mathrm{G}'\) a graded homomorphism of degree 0, such that \(v(\mathrm{F}^p \mathrm{G}) \subset \mathrm{F}^p \mathrm{G}'\) for all \(p\) and that the induced homomorphisms \(\bar{v}^{p,n}: (\mathrm{F}^p \mathrm{G})^n / (\mathrm{F}^{p+1} \mathrm{G})^n \to (\mathrm{F}^p \mathrm{G}')^n / (\mathrm{F}^{p+1} \mathrm{G}')^n\) satisfy \(\bar{v}^{p,p+q} \circ \beta^{p,q} = \beta'^{p,q} \circ u_{\infty}^{p,q}\) for every pair \((p, q)\) . Show that if \(u_2^{pq}\) is an isomorphism for every pair \((p, q)\) , then \(v\) is an isomorphism.

\begin{enumerate}
\def\labelenumi{\arabic{enumi})}
\setcounter{enumi}{14}
\tightlist
\item
  Let E be a spectral sequence of A-modules, which approaches a graded A-module G (exercise 14).
\end{enumerate}

\begin{enumerate}
\def\labelenumi{\alph{enumi})}
\tightlist
\item
  Assume \(E_{2}^{p,q}=0\) for p\textless0. Show that the maps \(\alpha_{r}\) induce for every q injective A-homomorphisms:
\end{enumerate}

\[
\mathrm{E} _ {\infty} ^ {0, q} \rightarrow \dots \rightarrow \mathrm{E} _ {3} ^ {0, q} \rightarrow \mathrm{E} _ {2} ^ {0, q}
\]

whence, by composition with \((\beta^{0q})^{-1}\) and with the passage to the quotient map, an A-homomorphism \(\gamma^q: G^q \to E_2^{0,q}\) .

\begin{enumerate}
\def\labelenumi{\alph{enumi})}
\setcounter{enumi}{1}
\tightlist
\item
  Assume \(E_{2}^{pq} = 0\) for q \textless{} 0. Show that the maps \(\alpha_{r}^{-1}\) induce surjective homomorphisms:
\end{enumerate}

\[
\mathrm{E} _ {2} ^ {p, 0} \rightarrow \mathrm{E} _ {3} ^ {p, 0} \rightarrow \dots \rightarrow \mathrm{E} _ {\infty} ^ {p, 0}
\]

whence, by composition with \(\beta^{p,0}\) , an A-homomorphism \(\delta^p:E_2^{p,0}\to G^p\) c) We assume that the hypotheses of a) and b) are satisfied, that is, \(E_{2}^{pq} = 0\) if p \textless{} 0 or q \textless{} 0. Show that the sequence

\[
0 \rightarrow \mathrm{E} _ {2} ^ {1, 0} \xrightarrow {\delta^ {1}} \mathrm{G} ^ {1} \xrightarrow {\gamma^ {1}} \mathrm{E} _ {2} ^ {0, 1} \xrightarrow {d _ {2}} \mathrm{E} _ {2} ^ {2, 0} \xrightarrow {\delta^ {2}} \mathrm{G} ^ {2}
\]

is exact.

\begin{enumerate}
\def\labelenumi{\alph{enumi})}
\setcounter{enumi}{3}
\item
  Assume \(\mathbf{E}_2^{pq} = 0\) if \(p > 0\) or \(q > 0\) ; define an exact sequence analogous to that of \(c\) ).
\item
  One assumes that there exist integers \(r, d\) , with \(d \geqslant r \geqslant 2\) , such that \(\mathrm{E}_{r}^{pq} = 0\) if \(p \neq 0\) , \(d\) . Define an exact sequence:
\end{enumerate}

\[
\dots \rightarrow \mathrm{E} _ {r} ^ {d, n - d} \rightarrow \mathrm{G} ^ {n} \rightarrow \mathrm{E} _ {r} ^ {0, n} \rightarrow \mathrm{E} _ {r} ^ {d, n + 1 - d} \rightarrow \mathrm{G} ^ {n + 1} \rightarrow \dots .
\]

If likewise \(E_{r}^{pq} = 0\) for \(q \neq 0\) , d, show that there exists an exact sequence:

\[
\dots \rightarrow \mathrm{E} _ {r} ^ {n, 0} \rightarrow \mathrm{G} ^ {n} \rightarrow \mathrm{E} _ {r} ^ {n - d, d} \rightarrow \mathrm{E} _ {r} ^ {n + 1, 0} \rightarrow \dots .
\]

\begin{enumerate}
\def\labelenumi{\alph{enumi})}
\setcounter{enumi}{5}
\tightlist
\item
  Assume \(\mathbf{E}_2^{pq} = 0\) if \(q \neq 0\) (resp. \(p \neq 0\) ). Show that the homomorphism
\end{enumerate}

\[
\delta^ {n}: \mathrm{E} _ {2} ^ {n, 0} \rightarrow \mathrm{G} ^ {n} (\text { resp. } \gamma^ {n}: \mathrm{G} ^ {n} \rightarrow \mathrm{E} _ {2} ^ {0, n})
\]

is bijective for every n.

¶ 16) Let C be an A-complex, \((F^{p}C)_{p\in\mathbf{Z}}\) a decreasing sequence of subcomplexes of C.

\begin{enumerate}
\def\labelenumi{\alph{enumi})}
\tightlist
\item
  For \(p,q,r\in\mathbf{Z}\) , let \(M_{r}^{pq}\) the set of elements x of \((F^{p}C)^{p+q}\) such that \(dx\in F^{p+r}C\) ; for \(r\geqslant0\) , we set \(E_{r}^{pq}=M_{r}^{pq}/(d(M_{r-1}^{p-r+1,q+r-2})+M_{r}^{p+1,q-1})\) and we denote by \(d_{r}:E_{r}^{pq}\to E_{r}^{p+r,q-r+1}\) the homomorphism deduced from \(d\) by passing to quotients. Show that there exists for \(r\geqslant0\) an isomorphism \(\alpha_{r}:H(E_{r})\to E_{r+1}\) , so that the sequence \((E_{r},d_{r})_{r\geqslant2}\) , equipped with the maps \(\alpha_{r}\) , is a spectral sequence. We have
\end{enumerate}

\[
\mathrm{E} _ {0} ^ {p q} = (\mathrm{F} ^ {p} \mathrm{C}) ^ {p + q} / (\mathrm{F} ^ {p + 1} \mathrm{C}) ^ {p + q}.
\]

\begin{enumerate}
\def\labelenumi{\alph{enumi})}
\setcounter{enumi}{1}
\item
  We assume that \(\bigcup_{p} \mathbf{F}^{p} \mathbf{C} = \mathbf{C}\) and that for every \(n\) , there exists an integer \(p\) such that \((\mathbf{F}^{p} \mathbf{C})^{n} = 0\) . Show that the spectral sequence defined in \(a)\) converges (exercise 14) to the graded A-module \(\mathrm{H}(\mathbf{C})\) , equipped with the sequence of graded submodules \(\mathbf{F}^{p} \mathbf{H}(\mathbf{C}) = \operatorname{Im} \mathbf{H}(i_{p})\) where \(i_{p}: \mathbf{F}^{p} \mathbf{C} \to \mathbf{C}\) is the canonical injection.
\item
  If \(\mathbf{F}^0\bar{\mathbf{C}} = \mathbf{C}\) the A-module \(\mathbf{E}_r^{pq}\) is zero for \(p < 0\) ; if \((\mathbf{F}^p\mathbf{C})^n = 0\) for \(p > n\) , we have \(\mathbf{E}_r^{pq} = 0\) for \(q < 0\) . d) Let \(\mathbf{C}'\) be a second complex, \((\mathbf{F}^p\mathbf{C}')_{p\in \mathbb{Z}}\) a decreasing sequence of subcomplexes of \(\mathbf{C}'\) , \(\mathbf{E}'\) the associated spectral sequence. Let \(u:\mathbf{C}\to \mathbf{C}'\) be a morphism of complexes such that \(u(\mathbf{F}^p\mathbf{C})\subset \mathbf{F}^p\mathbf{C}'\) for all \(p\) . Show that \(u\) induces a morphism of spectral sequences \(\tilde{u}:\mathbf{E}\to \mathbf{E}'\) .
\item
  With the notations of \(d\) ), let \(v: \mathbf{C} \to \mathbf{C}'\) be a second morphism such that \(v(\mathbf{F}^p \mathbf{C}) \subset \mathbf{F}^p \mathbf{C}'\) for all \(p\) , and let \(h\) a homotopy relating \(u\) to \(v\) . Let \(k\) be an integer \(\geqslant 0\) such that \(h(\mathbf{F}^p \mathbf{C}) \subset \mathbf{F}^{p-k} \mathbf{C}'\) for all \(p\) . Show that \(\tilde{u}_r = \tilde{v}_r\) for \(r > k\) .
\end{enumerate}

\begin{enumerate}
\def\labelenumi{\arabic{enumi})}
\setcounter{enumi}{16}
\tightlist
\item
  Let C be a bicomplex (X, p. 174, exercise 11).
\end{enumerate}

\begin{enumerate}
\def\labelenumi{\alph{enumi})}
\tightlist
\item
  One defines two decreasing sequences \(({}^{\prime}\mathbf{F}^p\mathbf{C})_{p\in \mathbf{Z}}\) and \(({}^{\prime\prime}\mathbf{F}^q\mathbf{C})_{q\in \mathbf{Z}}\) of subcomplexes of the complex associated with C by setting
\end{enumerate}

\[
\left(^ {\prime} \mathrm{F} ^ {p} \mathrm{C}\right) ^ {n} = \bigoplus_ {i \geqslant p} \mathrm{C} ^ {i, n - i} \quad \left(^ {\prime \prime} \mathrm{F} ^ {q} \mathrm{C}\right) ^ {n} = \bigoplus_ {j \geqslant q} \mathrm{C} ^ {n - j, j}.
\]

Let ′E and ″E be the corresponding spectral sequences (exercise 16, a)). Show that for every pair \((p, q)\) the A-module \({}^{\prime}E_2^{pq}\) (resp. \({}^{\prime\prime}E_2^{pq}\) ) is isomorphic to \(H^{\prime\prime p}(H^{\prime q}(C))\) (resp. \(H^{\prime p}(H^{\prime\prime q}(C))\) ).

\begin{enumerate}
\def\labelenumi{\alph{enumi})}
\setcounter{enumi}{1}
\item
  Suppose there exists an integer \(n\) such that \(C^{pq} = 0\) for \(p \geqslant n\) or for \(q \leqslant n\) . Show that the spectral sequence ′E converges to H(C). If likewise there exists an integer \(m\) such that \(C^{pq} = 0\) for \(p \leqslant m\) or for \(q \geqslant m\) , the spectral sequence ″E converges to H(C).
\item
  We assume that \(C^{pq} = 0\) for p \textless{} 0 or q \textless{} 0. Show that the homomorphism \(\delta^{p}: E_{2}^{p0} \to H^{p}(C)\) (exercise 15) comes, by passing to homology, from the canonical injection of complexes ``Z\^{}\{0\}(C) → C''. Likewise the homomorphism ``E\_\{2\}\^{}\{p0\} → H\^{}\{p\}(C)'' comes from the injection ``Z\^{}\{0\}(C) → C''.
\item
  Let \(\tilde{C}\) a second bicomplex, ′E and ″E the two associated spectral sequences. Show that every morphism of bicomplexes \(u: C \to \tilde{C}\) induces morphisms of spectral sequences ′u : ′E → ′E and ″u : ″E → ″E. Show that if two morphisms \(u\) , \(v\) of \(C\) to \(\tilde{C}\) are homotopic, the associated morphisms ′u and ′v (resp. ″u and ″v) are equal (use exercise 16, e)).
\end{enumerate}

\begin{enumerate}
\def\labelenumi{\arabic{enumi})}
\setcounter{enumi}{17}
\tightlist
\item
  Let E be a spectral sequence (exercise 14), such that the A-module \(\mathbf{E}_2\) has finite length. Let \(\mathcal{C}\) be the set of isomorphism classes of A-modules of finite length; show that one has \(\chi_{\mathcal{G}}(\mathbf{E}_2) = \chi_{\mathcal{G}}(\mathbf{E}_r) = \chi_{\mathcal{G}}(\mathbf{E}_{\infty})\) for all \(r \geqslant 2\) (X, p. 41). If moreover E converges to a graded A-module G, the latter has finite length and one has \(\chi_{\mathcal{G}}(\mathbf{G}) = \chi_{\mathcal{G}}(\mathbf{E}_2)\) .
\end{enumerate}

¶ 19) One calls a simplicial scheme a pair (K, F), where K is a set and F a set of finite subsets of K, called faces of K, such that every subset of a face is a face and every point of K belongs to at least one face. A simplicial map from a simplicial scheme (K, F) to a simplicial scheme (L, G) is a map from K to L which carries every face of K to a face of L. a) Let n be an integer ≥ 0. One calls an n-simplex of the simplicial scheme K every sequence (α₀, \ldots, αₙ) of elements of K such that the αᵢ belong to a single face of K. One denotes by Sₙ(K, A) the free A-module generated by the n-simplexes for n ≥ 0, and one sets Sₙ(K, A) = 0 for n \textless{} 0. One defines an A-homomorphism \(d_{n}:\mathbf{S}_{n}(\mathbf{K},\mathbf{A})\to \mathbf{S}_{n - 1}(\mathbf{K},\mathbf{A})\) by setting \(d_{n} = 0\) for \(n\leqslant 0\) and \(d_{n}(\alpha_{0},\dots,\alpha_{n}) = \sum_{i = 0}^{n}(-1)^{i}(\alpha_{0},\dots,\theta_{i},\dots,\alpha_{n})\)

for all \(n\) -simplex \((\alpha_0, \ldots, \alpha_n)\) (the sign over a letter means that it is to be omitted). Show that \(d_{n-1} \circ d_n = 0\) for all \(n\) , so that \((S(K, A), d)\) is a complex of A-modules. One defines in the same way an A-complex \(C(K, A)\) by setting \(C^n(K, A) = \text{Hom}_Z(S_n(K, Z), A)\) , \(\delta^n = \text{Hom}(d_n, 1_A)\) for all \(n\) . For \(n \geqslant 0\) , one sometimes denotes \(H_n(K, A)\) (resp. \(H^n(K, A)\) ) the A-module \(H_n(S(K, A))\) (resp. \(H^n(C(K, A)))\) . b) Let K, L be two simplicial schemes, \(f: K \to L\) a simplicial map. Show that \(f\) induces morphisms of complexes \(S(f): S(K, A) \to S(L, A)\) and \(C(f): C(L, A) \to C(K, A)\) , whence A-homomorphisms \(f_*: H(K, A) \to H(L, A)\) and \(f^*: H'(L, A) \to H'(K, A)\) .

\begin{enumerate}
\def\labelenumi{\alph{enumi})}
\setcounter{enumi}{2}
\tightlist
\item
  Let \(g: \mathbf{K} \to \mathbf{L}\) a second simplicial map. We say that \(f\) and \(g\) are simplicially homotopic if for every face \(\sigma\) of \(\mathbf{K}\) , the set \(f(\sigma) \cup g(\sigma)\) is a face of \(\mathbf{L}\) . Show that the morphisms \(\mathbf{S}(f)\) and \(\mathbf{S}(g)\) (resp. \(\mathbf{C}(f)\) and \(\mathbf{C}(g)\) ) are then homotopic (define a homotopy \(h\) by setting
\end{enumerate}

\[
h (\alpha_ {0}, \dots , \alpha_ {n}) = \sum_ {i = 0} ^ {n} (- 1) ^ {i} (f (\alpha_ {0}), \dots , f (\alpha_ {i}), g (\alpha_ {i}), \dots , g (\alpha_ {n}))
\]

for every n-simplex \((\alpha_{0}, \ldots, \alpha_{n})\) .

\begin{enumerate}
\def\labelenumi{\alph{enumi})}
\setcounter{enumi}{3}
\tightlist
\item
  A simplicial scheme K is said to be conical if there exists a point s of K such that for every face σ of K, σ ∪ \{s\} be a face. Show that there then exist homotopisms from S(K, A) and from C(K, A) onto A. e) Let D(K, A) denote the subcomplex of S(K, A) generated by the simplices (α₀, \ldots, αₙ) for which two of the αᵢ are equal, and by the elements of the form (α₀, \ldots, αₙ) - ε(σ). (ασ(0), \ldots, ασ(n)) for every n-simplex (α₀, \ldots, αₙ) and every permutation σ of \{0, \ldots, n\}Show that D(K, A) is a subcomplex of S(K, A), and that the quotient map p : S(K, A) → S(K, A)/D(K, A) is a homotopy equivalence (a homotopy inverse is obtained by choosing a total order on K and associating to the class modulo D(K, A) of a simplex (α₀, \ldots..., αₙ) the simplex ε(σ). (ασ(0), \ldots, ασ(n)), where σ is the permutation such that ασ(0) \textless{} \ldots{} \textless{} ασ(n), if all the αᵢ are distinct, and 0 otherwise).
\end{enumerate}

*20) Let R be a complete discrete valuation ring (AC, VI, § 3, no. 6) with residue field k of characteristic p \textgreater{} 0. Assume that the field of fractions of R has characteristic 0. Let v denote the normalized valuation of R, m the maximal ideal of R, and R* the multiplicative group of units of R. For every integer \(n \geqslant 0\) Let R*(n) be the subgroup of R* consisting of the elements x such that \(v(x - 1) \geqslant n\) . a) Let r be an integer, \(r > v(p)/(p - 1)\) . Show that the series

\[
\exp (x) = \sum_ {i = 0} ^ {+ \infty} x ^ {i} / i!
\]

converges for \(x \in \mathfrak{m}^r\) and defines an isomorphism from the additive group \(\mathfrak{m}^r\) onto the group \(\mathbf{R}^*(r)\) . b) Suppose that \(k\) is a finite field with \(q\) elements. Let \(n\) be an integer \(\geqslant 1\) , \(\mu_n\) the subgroup of elements \(x\) of \(\mathbf{R}^*\) such that \(x^n = 1\) . Show that the groups \(\mu_n\) and \(\mathbf{R}^*/(\mathbf{R}^*)^n\) are finite and that one has

\[
\operatorname{Card} \left(\mathrm{R} ^ {*} / \mathrm{R} ^ {* n}\right) / \operatorname{Card} \left(\mu_ {n}\right) = q v (n).
\]

(Let \(\mathcal{C}\) the set of classes of finite \(\mathbf{Z}\) -modules, \(\varphi : \mathrm{K}(\mathcal{C}) \to \mathbf{Q}^*\) be the homomorphism defined by

\[
\varphi ([ \mathrm{G} ]) = \text { Card } (\mathrm{G}) \quad \text { for } \quad \mathrm{G} \in \mathscr {C},
\]

\(C^{(r)}\) the complex of Z-modules such that \(C_{p}^{(r)} = 0\) for \(p \neq 0, 1\) , \(C_{0}^{(r)} = C_{1}^{(r)} = R^{*}(r)\) , \(d_{1}(x) = x^{n}\) for \(x \in R^{*}(r)\) . Prove that \(\varphi(H(C^{(r)}))\) is independent of r, then use a).

\begin{enumerate}
\def\labelenumi{\arabic{enumi})}
\setcounter{enumi}{20}
\tightlist
\item
  Assume that the ring A is a commutative algebra over a commutative ring k; the A-module \(D_{k}(A)\) of k-derivations of A is then identified with the dual of the A-module \(\Omega_{A/k}^{1}\) , which defines on \(\Omega_{A/k}\) a structure of \(\Lambda_{\mathbf{A}}(\mathbf{D}_{k}(\mathbf{A}))\) left module (cf. III, p. 165). Prove the formula
\end{enumerate}

\[
\begin{array}{l} \left(\mathbf {X} _ {0} \wedge \dots \wedge \mathbf {X} _ {p}\right) \rfloor d \omega = \sum_ {i = 0} ^ {p} (- 1) ^ {i} \mathbf {X} _ {i} \left(\left(\mathbf {X} _ {0} \wedge \dots \wedge \hat {\mathbf {X}} _ {i} \wedge \dots \wedge \mathbf {X} _ {p}\right) \rfloor \omega\right) \\ \quad + \sum_ {0 \leqslant i <   j \leqslant p} (- 1) ^ {i + j} \left([ \mathbf {X} _ {i}, \mathbf {X} _ {j} ] \wedge \mathbf {X} _ {0} \wedge \dots \wedge \hat {\mathbf {X}} _ {i} \wedge \dots \wedge \hat {\mathbf {X}} _ {r} \wedge \dots \wedge \mathbf {X} _ {p}\right) \rfloor d \omega . \end{array}
\]

where \(\omega\in\Omega_{\mathbf{A}/k}^{p},\mathbf{X}_{0},\ldots,\mathbf{X}_{p}\in\mathbf{D}_{k}(\mathbf{A})\) and where the sign \(^{\wedge}\) on a letter means that it is to be omitted. (Reduce to the case \(p=1\) .)

\begin{enumerate}
\def\labelenumi{\arabic{enumi})}
\setcounter{enumi}{21}
\tightlist
\item
  We keep the conventions of the preceding exercise; if M is an A-module, \(\nabla\) a connection on M and if \(X \in D_k(A)\) , we denote by \(\nabla_X\) the \(k\) -endomorphism \((1_M \otimes X) \circ \nabla\) of M.
\end{enumerate}

\begin{enumerate}
\def\labelenumi{\alph{enumi})}
\item
  Let \(\mathbf{M}_1, \mathbf{M}_2\) two A-modules, equipped with connections \(^1\nabla\) , \(^2\nabla\) ; show that one defines a connection \(\nabla\) over \(\mathbf{M}_1 \oplus \mathbf{M}_2\) (resp. \(\mathbf{M}_1 \otimes_{\mathbf{A}} \mathbf{M}_2\) , resp. \(\operatorname{Hom}_{\mathbf{A}}(\mathbf{M}_1, \mathbf{M}_2)\) ) by setting \(\nabla(m_1 + m_2) = ^1\nabla(m_1) + ^2\nabla(m_2)\) for \(m_1 \in \mathbf{M}_1\) , \(m_2 \in \mathbf{M}_2\) (resp. \(\nabla_{\mathbf{X}}(m_1 \otimes m_2)) = ^1\nabla_{\mathbf{X}}(m_1) \otimes m_2 + m_1 \otimes ^2\nabla_{\mathbf{X}}(m_2)\) for \(\mathbf{X} \in \mathbf{D}_k(\mathbf{A})\) , \(m_1 \in \mathbf{M}_1\) and \(m_2 \in \mathbf{M}_2\) , resp. \(\nabla_{\mathbf{X}}(u) = ^2\nabla_{\mathbf{X}} \circ u - u \circ ^1\nabla_{\mathbf{X}}\) for \(\mathbf{X} \in \mathbf{D}_k(\mathbf{A})\) , \(u \in \operatorname{Hom}_{\mathbf{A}}(\mathbf{M}_1, \mathbf{M}_2)\) .
\item
  Let M be an A-module, \(\nabla: \mathbf{M} \to \mathbf{M} \otimes_{\mathbf{A}} \Omega_{\mathbf{A}/k}^{1}\) a connection on M. If X, Y ∈ D \(_{k}\) (A), we denote by R(X, Y) the A-endomorphism of M such that R(X, Y)(m) = \textless R(m), X ∧ Y\textgreater{} for all m ∈ M. Prove the formula: R(X, Y) = {[}∇X, ∇Y{]} - ∇{[}X,Y{]} (use Exercise 21).
\item
  Assume the A-module M is projective, so that the curvature homomorphism R is identified with an element of \(\operatorname{End}_{\mathbf{A}}(\mathbf{M}) \otimes_{\mathbf{A}} \Omega_{\mathbf{A}/k}^{2}\) . If one equips \(\operatorname{End}_{\mathbf{A}}(\mathbf{M})\) with the connection \(\nabla\) defined in \(a\) ), show that one has \(\nabla^{2}\mathbf{R} = 0\) .
\end{enumerate}

\exercisesection{Resolutions}

\begin{enumerate}
\def\labelenumi{\arabic{enumi})}
\tightlist
\item
  Let \(a, b\) two elements of A, such that the left annihilator of \(a\) (resp. \(b\) ) is the ideal \(Ab\) (resp. \(Aa\) ). Show that the sequence
\end{enumerate}

\[
\dots \longrightarrow \mathrm{A} _ {s} \xrightarrow {\delta_ {a}} \mathrm{A} _ {s} \xrightarrow {\delta_ {b}} \mathrm{A} _ {s} \xrightarrow {\delta_ {a}} \mathrm{A} _ {s} \longrightarrow \dots \longrightarrow \mathrm{A} _ {s} \xrightarrow {\delta_ {a}} \mathrm{A} _ {s} \longrightarrow \mathrm{A} / \mathrm{A} a \longrightarrow 0,
\]

where \(\delta_{a}\) (resp. \(\delta_{b}\) ) denotes right multiplication by \(a\) (resp. \(b\) ), defines a free resolution of the A-module A/Aa. Show that this construction applies in the following two cases:

\(\alpha)\) Let \(\mathbf{A}_0\) a ring; we take \(\mathrm{A} = \mathrm{A}_0(\varepsilon)\) (X, p. 27), \(a = b = \varepsilon\) .

β) Let G be a finite cyclic group, σ a generator of G, k a commutative ring. We take

\[
\mathbf {A} = k ^ {(\mathbf {G})}, \quad a = 1 - e _ {\sigma}, \quad b = \sum_ {g \in \mathbf {G}} e _ {g}.
\]

\begin{enumerate}
\def\labelenumi{\arabic{enumi})}
\setcounter{enumi}{1}
\tightlist
\item
  We assume the ring A is commutative. Let B be an A-algebra; for \(n \geqslant 0\) , we denote by \(\mathcal{A}^{n}(\mathbf{B})\) the tensor product over A of \((n + 1)\) modules equal to B. We set \(\mathcal{A}^{n}(\mathbf{B}) = 0\) for \(n < 0\) and \(\mathcal{A}(\mathbf{B}) = \bigoplus \mathcal{A}^{n}(\mathbf{B})\) We define a graded A-endomorphism of ascending degree +1 of \(\mathcal{A}(\bar{\mathbf{B}})\) by setting:
\end{enumerate}

\[
d (b _ {0} \otimes \dots \otimes b _ {n}) = \sum_ {i = 0} ^ {n + 1} (- 1) ^ {i} b _ {0} \otimes \dots \otimes b _ {i - 1} \otimes 1 _ {\mathbf {B}} \otimes b _ {i} \otimes \dots \otimes b _ {n} \quad \text { for } b _ {0}, \dots , b _ {n} \in \mathbf {B}.
\]

\begin{enumerate}
\def\labelenumi{\alph{enumi})}
\tightlist
\item
  Show that \(d \circ d = 0\) , so that \((\mathcal{A}(\mathbf{B}), d)\) is a complex of A-modules. If M is an A-module, we denote by \(\mathcal{A}(\mathbf{B}, \mathbf{M})\) the complex such that \(\mathcal{A}^{n}(\mathbf{B}, \mathbf{M}) = \mathcal{A}^{n}(\mathbf{B}) \otimes_{\mathbf{A}} \mathbf{M}\) , with differential \(d \otimes 1_{\mathbf{M}}\) . b) Suppose that the algebra B is a faithfully flat A-module (X, p. 169, exercise 9). Show that the complex \(\mathcal{A}(\mathbf{B}, \mathbf{M})\) defines a right resolution of the A-module M. (It suffices to verify that the complex \(\mathbf{B} \otimes_{\mathbf{A}} \mathcal{A}(\mathbf{B}, \mathbf{M})\) defines a resolution of \(\mathbf{B} \otimes_{\mathbf{A}} \mathbf{M}\) ; construct an A-linear homotopism between these two complexes.)
\end{enumerate}

\begin{enumerate}
\def\labelenumi{\arabic{enumi})}
\setcounter{enumi}{2}
\tightlist
\item
  Let B be a ring, a a two-sided ideal of B, b a left ideal of B containing a. Denote by A the ring B/a, by \(p: \mathbf{B} \to \mathbf{A}\) the quotient homomorphism and by \(b'\) the ideal \(p(\mathbf{b}) \subset \mathbf{A}\) . a) Consider the sequence of canonical injections:
\end{enumerate}

\[
\dots \rightarrow a ^ {n} b \rightarrow a ^ {n} \rightarrow a ^ {n - 1} b \rightarrow a ^ {n - 1} \rightarrow \dots \rightarrow a \rightarrow b \rightarrow B.
\]

Show that the sequence obtained by tensor product with B/a

\[
\dots \rightarrow \mathfrak {a} ^ {n} \mathfrak {b} / \mathfrak {a} ^ {n + 1} \mathfrak {b} \rightarrow \mathfrak {a} ^ {n} / \mathfrak {a} ^ {n + 1} \rightarrow \dots \rightarrow \mathfrak {a} / \mathfrak {a} ^ {2} \rightarrow \mathfrak {b} / \mathfrak {a b} \rightarrow \mathbf {A}
\]

defines a left resolution of the A-module A/b'.

\begin{enumerate}
\def\labelenumi{\alph{enumi})}
\setcounter{enumi}{1}
\item
  Suppose that the ideals a and b are free left A-modules. Show that the preceding resolution is a free resolution of the A-module A/b'.
\item
  Let G be a group, k a commutative ring, \(A = \dot{k}^{(G)}\) the algebra of G over k, \(\varepsilon : A \to k\) the homomorphism of k-algebras such that \(\varepsilon(e_g) = 1\) for all \(g \in G\) . One endows k with the structure of A-module (on the left) associated with the homomorphism \(\varepsilon\) . Let \((g_i)_{i \in I}\) be a generating family of G, F the free group constructed on I, \(B = k^{(F)}\) , \(\pi : F \to G\) the homomorphism such that \(\pi(i) = g_i\) for \(i \in I\) , \(p : B \to A\) the homomorphism of k-algebras deduced from \(\pi\) . Show that the ideals a = Ker (p) and b = Ker ( \(\varepsilon \circ p\) ) are free left B-modules; deduce from this a free resolution of the \(k^{(G)}\) -module k.
\end{enumerate}

(According to I, p. 147, exercise 20, the group Ker (π) admits a basic family (rα)α∈J; show that the elements (eᵣα - 1), for α ∈ J (resp. (eᵢ - 1), for i ∈ I) form a basis of the left B-module a (resp. b).)

\begin{enumerate}
\def\labelenumi{\alph{enumi})}
\setcounter{enumi}{3}
\tightlist
\item
  With the notations of c), we set \(\mathbf{R} = \operatorname{Ker}(\pi)\) ; let \(\varphi: \mathbf{R} \to \mathbf{R}/(\mathbf{R}, \mathbf{R})\) be the quotient map. We endow the \(k\) -module \(k \otimes_{\mathbf{Z}} \mathbf{R}/(\mathbf{R}, \mathbf{R})\) of the A-module structure such that
\end{enumerate}

\[
e _ {\pi (f)} \left(1 \otimes \varphi (r)\right) = 1 \otimes \varphi (f r f ^ {- 1})
\]

for all \(r \in \mathbb{R}\) , \(f \in \mathbb{F}\) . Show that there exists an exact sequence of A-modules:

\[
0 \rightarrow k \otimes_ {\mathbf {Z}} \mathrm{R} / (\mathrm{R}, \mathrm{R}) \rightarrow \mathrm{A} ^ {(I)} \rightarrow \mathrm{A} \rightarrow k \rightarrow 0.
\]

\begin{enumerate}
\def\labelenumi{\arabic{enumi})}
\setcounter{enumi}{3}
\tightlist
\item
  Assume the ring A is commutative.
\end{enumerate}

Let G be a group, and K a simplicial scheme (X, p. 177, exercise 19). An action of G on K is a homomorphism from G into the group of bijective simplicial maps from K to itself. a) An action of G on K defines an action of G on the A-module S(K, A) (loc. cit.); show that this action makes S(K, A) into a complex of A \(^{(G)}\) -modules.

\begin{enumerate}
\def\labelenumi{\alph{enumi})}
\setcounter{enumi}{1}
\item
  Assume that G acts freely on the underlying set of K. Show that the A \(^{(G)}\) -modules S \(_{n}\) (K, A) are free.
\item
  Assume that the simplicial scheme K is conical. Show that the \(\mathbf{A}^{(\mathbf{G})}\) -complex \(\mathbf{S}(\mathbf{K},\mathbf{A})\) defines a left resolution of A, equipped with the structure of \(\mathbf{A}^{(\mathbf{G})}\) -module for which \(e_g a = a\) for all \(g \in \mathbf{G}\) , \(a \in \mathbf{A}\) .
\item
  Let \(|\mathbf{G}|\) the simplicial scheme \((\mathbf{G},\mathcal{F})\) , where \(\mathcal{F}\) is the set of finite subsets of G. The group G acts by left translations on the simplicial scheme \(|\mathbf{G}|\) ; show that the \(\mathbf{A}^{(\mathbf{G})}\) -complex S(\textbar G\textbar, A) is a free resolution of the \(\mathbf{A}^{(\mathbf{G})}\) -module A.
\item
  Let \(\mathbf{B}(\mathbf{A}^{(G)},\mathbf{A})\) the standard resolution of the \(\mathbf{A}^{(G)}\) -module A; for \(n\geqslant 0\) , the \(\mathbf{A}^{(G)}\) -module \(\mathbf{B}_n(\mathbf{A}^{(G)},\mathbf{A})\) is identified with \((\mathbf{A}^{(G)})^{\otimes (n + 1)}\) . Show that the A-linear homomorphism \(u:\mathrm{S}(|\mathrm{G}|,\mathrm{A})\to \mathrm{B}(\mathrm{A}^{(G)},\mathrm{A})\) , defined by
\end{enumerate}

\[
u \left(g _ {0}, \dots , g _ {n}\right) = g _ {0} \otimes g _ {0} ^ {- 1} g _ {1} \otimes \dots \otimes g _ {n - 1} ^ {- 1} g _ {n} \quad \text { for } g _ {0}, \dots , g _ {n} \text { in } G,
\]

is an isomorphism of \(A^{(G)}\) -complexes.

\begin{enumerate}
\def\labelenumi{\arabic{enumi})}
\setcounter{enumi}{4}
\tightlist
\item
  \begin{enumerate}
  \def\labelenumii{\alph{enumii})}
  \tightlist
  \item
    Let \(0 \to K \to P \to M \to 0\) , \(0 \to K' \to P' \to M \to 0\) two exact sequences of A-modules, with \(P\) and \(P'\) projective. Show that there exists a commutative diagram:
  \end{enumerate}
\end{enumerate}

\[
\begin{array}{c c c}0 \rightarrow \mathrm{K} \oplus \mathrm{P} ^ {\prime} \rightarrow \mathrm{P} \oplus \mathrm{P} ^ {\prime} \rightarrow \mathrm{M} \rightarrow 0\\\downarrow^ {v}&\downarrow^ {u}&\downarrow^ {l _ {M}}\\0 \rightarrow \mathrm{P} \oplus \mathrm{K} ^ {\prime} \rightarrow \mathrm{P} \oplus \mathrm{P} ^ {\prime} \rightarrow \mathrm{M} \rightarrow 0\end{array}
\]

where the horizontal lines are exact, and where \(u\) and \(v\) are isomorphisms (cf. II, p. 183, exercise 4). b) Let \(n\) be an integer \(\geqslant 0\) , \((\mathbf{P}, p)\) and \((\mathbf{P}', p')\) be two left resolutions of the A-module M, such that \(\mathbf{P}_j = \mathbf{P}_j' = 0\) for \(j > n\) and that the A-modules \(\mathbf{P}_i\) and \(\mathbf{P}_i'\) be projective for \(0 \leqslant i \leqslant n - 1\) . Show that the A-modules \(\bigoplus_{i \geqslant 0} (\mathbf{P}_{2i} \oplus \mathbf{P}_{2i+1}')\) and \(\bigoplus_{i \geqslant 0} (\mathbf{P}_{2j}^{\prime} \oplus \mathbf{P}_{2j+1})\) are isomorphic.

\begin{enumerate}
\def\labelenumi{\alph{enumi})}
\setcounter{enumi}{2}
\tightlist
\item
  Under the assumptions of \(b\) ), show that there exist two projective complexes \(\mathbf{R}\) , \(\mathbf{R}'\) of zero homology, such that \(\mathbf{R}_i = \mathbf{R}_i' = 0\) for \(i \notin [0, n]\) , and an isomorphism of resolutions \(u: \mathbf{P} \oplus \mathbf{R} \to \mathbf{P}' \oplus \mathbf{R}'\) .
\end{enumerate}

¶6). Let M be a left A-module. One calls a presentation of length n, or an n-presentation of M, an exact sequence

\[
\mathrm{L} _ {n} \rightarrow \mathrm{L} _ {n - 1} \rightarrow \dots \rightarrow \mathrm{L} _ {1} \rightarrow \mathrm{L} _ {0} \rightarrow \mathrm{M} \rightarrow 0
\]

where \(L_{i}\) is a free left A-module ( \(0 \leqslant i \leqslant n\) ). The presentation is said to be finite if all the \(L_{i}\) are finitely generated free modules.

If M is a finitely generated left A-module, one denotes by \(\lambda(M)\) the least upper bound (finite or equal to \(+\infty\) ) of the integers \(n \geqslant 0\) such that M admits a finite \(n\) -presentation. If M is not finitely generated, one sets \(\lambda(M) = -1\) .

\begin{enumerate}
\def\labelenumi{\alph{enumi})}
\item
  Let \(0 \to \mathbf{P} \to \mathbf{N} \to \mathbf{M} \to 0\) an exact sequence of left A-modules. One then has \(\lambda(\mathbf{N}) \geqslant \inf (\lambda(\mathbf{P}), \lambda(\mathbf{M}))\) .
\item
  Let \(L_{n} \xrightarrow{d_{n}} L_{n-1} \to \cdots \to L_{0} \to M \to 0\) a finite n-presentation of M; if \(\lambda(M) > n\) show that Ker \((d_{n})\) is a finitely generated A-module (use exercise 5, b)).
\item
  Show that, under the hypotheses of a), one has
\end{enumerate}

\[
\lambda (\mathbf {M}) \geqslant \inf \left(\lambda (\mathbf {N}), \lambda (\mathbf {P}) + 1\right).
\]

\((\mathrm{If} n \leqslant \inf (\lambda(\mathbf{N}), \lambda(\mathbf{P}) + 1)\) , show by induction on \(n\) that \(\lambda(\mathbf{M}) \geqslant n\) reasoning as in \(a\) and using \(b)\) .)

\begin{enumerate}
\def\labelenumi{\alph{enumi})}
\setcounter{enumi}{3}
\tightlist
\item
  Show that, under the hypotheses of a), one has
\end{enumerate}

\[
\lambda (\mathbf {P}) \geqslant \inf \left(\lambda (\mathbf {N}), \lambda (\mathbf {M}) - 1\right).
\]

(Same method as in c).) Deduce that, if \(\lambda(\mathbf{N}) = +\infty\) , then \(\lambda(\mathbf{M}) = \lambda(\mathbf{P}) + 1\) .

\begin{enumerate}
\def\labelenumi{\alph{enumi})}
\setcounter{enumi}{4}
\tightlist
\item
  Deduce from a), c) and d) that, if \(\mathbf{N} = \mathbf{M} \oplus \mathbf{P}\) , we have
\end{enumerate}

\[
\lambda (\mathbf {N}) = \inf \left(\lambda (\mathbf {M}), \lambda (\mathbf {P})\right).
\]

In particular, for N to admit a finite presentation it is necessary and sufficient that the same be true for M and P. f) Let \(N_{1}\) , \(N_{2}\) be two submodules of an A-module M. Suppose that \(N_{1}\) and \(N_{2}\) admit a finite presentation. For \(N_{1} + N_{2}\) to admit a finite presentation, it is necessary and sufficient that \(N_{1} \cap N_{2}\) be finitely generated.

\begin{enumerate}
\def\labelenumi{\arabic{enumi})}
\setcounter{enumi}{6}
\tightlist
\item
  \begin{enumerate}
  \def\labelenumii{\alph{enumii})}
  \tightlist
  \item
    With the notations of exercise 6, show that, if M is a projective module, one has
  \end{enumerate}
\end{enumerate}

\[
\lambda (\mathbf {M}) = - 1 \quad \text { or } \quad \lambda (\mathbf {M}) = + \infty .
\]

If A is a left Noetherian ring, then, for every A-module M, one has \(\lambda(M) = -1\) or \(\lambda(M) = +\infty\) . b) If a is a left ideal of a ring A, which is not finitely generated, \(A_s/a\) is a monogenic A-module which does not admit a finite presentation, in other words \(\lambda(A_s/a) = 0\) (cf. X, p. 11, prop. 6).

\begin{enumerate}
\def\labelenumi{\alph{enumi})}
\setcounter{enumi}{2}
\item
  Give an example of a monogenic left ideal a of a ring A such that \(A_{s}/a\) (which is of finite presentation) admits a dual which is not a finitely generated right A-module.
\item
  Let K be a commutative field, E the vector space \(\mathbf{K}^{(\mathbf{N})}\) , \((e_n)\) the canonical basis of E, T the tensor algebra of E, of which a basis is therefore formed by the finite products \(e_{i_1} e_{i_2} \ldots e_{i_k}\) ( \(k \geqslant 0\) , \(i_j \in \mathbf{N}\) for all \(j\) ). For a given integer \(n\) let b be the two-sided ideal of T generated by the products \(e_1 e_0, e_2 e_1, \ldots, e_n e_{n-1}\) and \(e_{n+k} e_n\) for all \(k \geqslant 1\) ; let A be the quotient ring T/b, and for every integer \(m\) , let \(a_{m}\) the canonical image of \(e_{m}\) in A. Show that, if \(\mathbf{M} = \mathbf{A}_{s}/\mathbf{A}a_{0}\) , we have \(\lambda(\mathbf{M}) = n\) (observe that, for \(m \leqslant n - 1\) the left annihilator of \(a_{m}\) is \(\mathbf{A}a_{m+1}\) and use exercise 6, b)).
\end{enumerate}

\begin{enumerate}
\def\labelenumi{\arabic{enumi})}
\setcounter{enumi}{7}
\tightlist
\item
  Let C be a commutative ring, E, F two C-modules. Show that one has \(\lambda(E \otimes_{C} F) \geqslant \inf(\lambda(E), \lambda(F))\) .
\item
  Let \(\rho: A \to B\) a ring homomorphism, M an A-module.
\end{enumerate}

\begin{enumerate}
\def\labelenumi{\alph{enumi})}
\item
  If B is a flat right A-module and if M admits a finite n-presentation, the B-module B \(\otimes_{\mathbf{A}}\) M admits a finite n-presentation.
\item
  If B is a faithfully flat right A-module (X, p. 169, exercise 9), and if the B-module B ⊗\_A M admits a finite n-presentation, then M admits a finite n-presentation (use exercise 6, b) and exercise 12, p. 169).
\end{enumerate}

\begin{enumerate}
\def\labelenumi{\arabic{enumi})}
\setcounter{enumi}{9}
\tightlist
\item
  Let E be an A-module. One says that E is pseudo-coherent if every finitely generated submodule of E is of finite presentation; every submodule of a pseudo-coherent module is pseudo-coherent. One says that E is coherent if it is pseudo-coherent and finitely generated (hence of finite presentation).
\end{enumerate}

\begin{enumerate}
\def\labelenumi{\alph{enumi})}
\item
  Let \(0 \to E' \to E \to E'' \to 0\) an exact sequence of A-modules. Show that, if E is pseudo-coherent (resp. coherent) and \(E'\) finitely generated, \(E''\) is pseudo-coherent (resp. coherent). Show that, if \(E'\) and \(E''\) are pseudo-coherent (resp. coherent), then so is E. Show that, if E and \(E''\) are coherent, so is \(E'\) (use X, p. 11, prop. 6).
\item
  Let E be a coherent A-module, E' a pseudo-coherent (resp. coherent) A-module. Show that, for every homomorphism \(u: \mathrm{E} \to \mathrm{E}'\) , \(\operatorname{Im}(u)\) and \(\operatorname{Ker}(u)\) are coherent and Coker ( \(u\) ) pseudo-coherent (resp. coherent) (use \(a\) )).
\item
  Show that every direct sum (resp. every finite direct sum) of pseudo-coherent (resp. coherent) modules is a pseudo-coherent (resp. coherent) module.
\item
  If E is a pseudo-coherent module and if M, N are coherent submodules of E, show that \(M + N\) and \(M \cap N\) are coherent (use a) and c)).
\item
  Suppose A is commutative. Show that, if E is a coherent A-module and F a coherent (resp. pseudo-coherent) A-module, Hom\_A(E, F) is a coherent (resp. pseudo-coherent) A-module. (Reduce to the case where F is coherent; consider a finite presentation of E, then use b).)
\end{enumerate}

¶11) a) Let A be a ring. Show that the following properties are equivalent:

\(\alpha)\) The A-module \(\mathbf{A}_s\) is coherent (exercise 10).

β) Every A-module of finite presentation is coherent.

γ) For every finitely presented right A-module M, the dual M* is a finitely generated A-module.

δ) For every set I, the right A-module \(A_{d}^{I}\) is flat.

ε) Any product of flat right A-modules is flat.

(To prove that \(\alpha\) ) implies \(\beta\) ), use Exercise 10, b). To see that \(\delta\) ) implies \(\gamma\) ), use X, p. 14, th. 1; to show that \(\alpha\) ) implies \(\varepsilon\) ), use X, p. 170, exercise 18.)

A ring satisfying these properties is said to be left coherent (or simply coherent), and the notion of a right coherent ring is defined similarly.

\begin{enumerate}
\def\labelenumi{\alph{enumi})}
\setcounter{enumi}{1}
\item
  Show that a left Noetherian ring is coherent. Give an example of a right Artinian ring that is not coherent (cf. VIII, § 1, exercise 4).
\item
  Show that an absolutely flat ring (X, p. 169, exercise 16) is coherent.
\item
  Show that, if A is a coherent ring and M is an A-module, we have \(\lambda(M) = -1\) or \(\lambda(M) = 0\) or \(\lambda(M) = +\infty\) (exercise 6). In particular, every finitely presented A-module admits a resolution by finitely generated free modules.
\item
  Let \((\mathbf{A}_{\alpha}, \varphi_{\beta \alpha})\) an inductive system of rings whose index set is filtered, and let \(\mathbf{A} = \varinjlim \mathbf{A}_{\alpha}\) . We assume that for \(\alpha \leqslant \beta\) , \(\mathbf{A}_{\beta}\) is a flat right \(\mathbf{A}_{\alpha}\) -module. Show that if the \(\mathbf{A}_{\alpha}\) are coherent, so is \(\mathbf{A}\) (Observe that \(\mathbf{A}\) is a flat right \(\mathbf{A}_{\alpha}\) -module for all \(\alpha\) and that for every finitely generated left ideal \(\alpha \subset \mathbf{A}\) there exists an index \(\alpha\) and an ideal \(\alpha_{\alpha}\) of \(\mathbf{A}\) such that \(\alpha\) is isomorphic to \(\mathbf{A} \otimes_{\mathbf{A}} a_{\alpha'}\) .)
\item
  Deduce from e) that every polynomial ring (for any finite or infinite set of indeterminates) over a commutative noetherian ring is coherent. Deduce from this that a quotient ring of a coherent ring is not necessarily coherent.
\item
  For A to be left coherent, it is necessary and sufficient that the left annihilator of every element of A be finitely generated, and that the intersection of two finitely generated left ideals be finitely generated (use exercise 6, f)).
\end{enumerate}

\begin{enumerate}
\def\labelenumi{\arabic{enumi})}
\setcounter{enumi}{11}
\item
  Let c be an infinite cardinal. Suppose that every left ideal of the ring A admits a generating set of cardinality ≤ c. Show that every A-module generated by a set of elements of cardinality ≤ c admits a left resolution by free A-modules of rank ≤ c (cf. VIII, § 1, exercise 13).
\item
  Let M be an A-module. A minimal injective resolution of M is an injective resolution \((e, 1)\) of M such that \(I^n\) is an injective envelope of \(Z^n(I)\) for all \(n \geqslant 0\) .
\end{enumerate}

\begin{enumerate}
\def\labelenumi{\alph{enumi})}
\tightlist
\item
  Every A-module admits minimal injective resolutions; any two such resolutions are isomorphic.
\item
  Let I, I' be two injective resolutions of M, with I minimal; let \(f: I \to I'\) and \(g: I' \to I\) two morphisms of resolutions. Then f is injective, g is surjective, and it is the direct sum of the subcomplexes Im(f) (isomorphic to I) and Ker(g) (with zero homology).
\end{enumerate}

\begin{enumerate}
\def\labelenumi{\arabic{enumi})}
\setcounter{enumi}{13}
\tightlist
\item
  Assume that the ring A is local Noetherian; denote its maximal ideal by m. Let L be a complex of finitely generated A-modules, zero in degrees \textless{} 0, M a finitely generated A-module, and q : L → M a morphism. We make the following hypotheses:
\end{enumerate}

\begin{enumerate}
\def\labelenumi{(\roman{enumi})}
\item
  The homomorphism \(1 \otimes q : (\mathbf{A}/\mathfrak{m}) \otimes_{\mathbf{A}} \mathbf{L}_{0} \to (\mathbf{A}/\mathfrak{m}) \otimes_{\mathbf{A}} \mathbf{M}\) is injective.
\item
  One has \(d_{i}(\mathbf{L}_{i})\subset \mathfrak{m}\mathbf{L}_{i - 1}\) for \(i\geqslant 1\) , and the map
\end{enumerate}

\[
\bar {d} _ {i}: \mathrm{L} _ {i} / \mathrm{mL} _ {i} \rightarrow \mathrm{mL} _ {i - 1} / \mathrm{m} ^ {2} \mathrm{L} _ {i - 1} \quad \text { is injective }.
\]

Let (P, p) be a projective resolution of M, \(u : L \to P\) a morphism such that \(p \circ u = q\) Show that u is injective and that \(u(L)\) is a direct summand of P as a graded A-module (reduce to the case where the resolution P is minimal, and use VIII, § 8, no. 3, cor. 3).

15) Let A be a Noetherian local ring, m its maximal ideal, k = A/m. Let P be a minimal projective resolution of the A-module k; for \(n \geqslant 0\) , we denote by \(b_{n}\) the rank of the free A-module \(P_{n}\) .

\begin{enumerate}
\def\labelenumi{\alph{enumi})}
\item
  Show that we have \(b_0 = 1\) and that \(b_1 = \dim_k (m/m^2)\) is the minimal number of generators of \(m\) .
\item
  If A is commutative, we have \(b_{2} \geqslant \frac{1}{2} b_{1}(b_{1} - 1)\) (consider the exact sequence \(\mathbf{P}_{2} \to m\mathbf{P}_{1} \to m^{2} \to 0\) ). c) Assume the ring A is Artinian (not necessarily commutative). Prove the inequality \(b_{2} \geqslant \frac{1}{4} b_{1}^{2}\) (let \(\mathbf{M}_{l}\) the submodule of \(\mathbf{P}_{2}\) of elements \(x\) such that \(d_{2}(x) \in m^{l}\mathbf{P}_{1}\) ; show that we have
\end{enumerate}

\[
\mathbf {P} _ {2} \subset \mathbf {M} _ {l + 1}
\]

and that there exists an exact sequence \(0 \to \mathbf{M}_l / \mathbf{M}_{l+1} \to \mathfrak{m}^l \mathbf{P}_1 / \mathfrak{m}^{l+1} \mathbf{P}_1 \to \mathfrak{m}^{l+1} / \mathfrak{m}^{l+2} \to 0\) Deduce from this that the polynomial \(p(t) = \sum_{l} (\dim_k (\mathfrak{m}^l / \mathfrak{m}^{l+1})) t^l\) satisfies \((b_2 t^2 - b_1 t + 1)p(t) \geqslant 0\) for \(0 \leqslant t \leqslant 1\) . Conclude

by examining the cases separately \(b_{1} = 1, b_{1} \geqslant 2\) .

\begin{enumerate}
\def\labelenumi{\arabic{enumi})}
\setcounter{enumi}{15}
\tightlist
\item
  \begin{enumerate}
  \def\labelenumii{\alph{enumii})}
  \tightlist
  \item
    Let \(0 \to \mathbf{M}' \to \mathbf{M} \to \mathbf{M}'' \to 0\) an exact sequence of A-modules, \((\mathbf{P}', p')\) (resp. \((\mathbf{P}'', p'')\) ) a projective resolution of \(\mathbf{M}'\) (resp. \(\mathbf{M}''\) ). Show that there exists a projective resolution \((\mathbf{P}, p)\) of \(\mathbf{M}\) and a commutative diagram:
  \end{enumerate}
\end{enumerate}

\[
\begin{array}{c} 0 \to \mathrm{P} ^ {\prime} \to \mathrm{P} \to \mathrm{P} ^ {\prime \prime} \to 0 \\ \downarrow^ {p ^ {\prime}} \quad \downarrow^ {p} \quad \downarrow^ {p ^ {\prime \prime}} \\ 0 \to \mathrm{M} ^ {\prime} \to \mathrm{M} \to \mathrm{M} ^ {\prime \prime} \to 0 \end{array}
\]

where the horizontal lines are exact.

\begin{enumerate}
\def\labelenumi{\alph{enumi})}
\setcounter{enumi}{1}
\tightlist
\item
  Let \(0 \to \mathbf{N}' \to \mathbf{N} \to \mathbf{N}'' \to 0\) a second exact sequence of A-modules, \((\mathbf{Q}, q)\) (resp. \((\mathbf{Q}', q')\) , resp. \((\mathbf{Q}'', q'')\) ) a projective resolution of \(\mathbf{N}\) (resp. \(\mathbf{N}'\) , resp. \(\mathbf{N}''\) ), so that there exists a commutative diagram with exact rows:
\end{enumerate}

\[
\begin{array}{c}0 \rightarrow Q ^ {\prime} \rightarrow Q \rightarrow Q ^ {\prime \prime} \rightarrow 0\\\downarrow^ {q ^ {\prime}} \quad \downarrow^ {q} \quad \downarrow^ {q ^ {\prime \prime}}\\0 \rightarrow N ^ {\prime} \rightarrow N \rightarrow N ^ {\prime \prime} \rightarrow 0\end{array}
\]

Let, on the other hand:

\[
\begin{array}{c} 0 \to \mathrm{M} ^ {\prime} \to \mathrm{M} \to \mathrm{M} ^ {\prime \prime} \to 0 \\ \downarrow^ {u ^ {\prime}} \quad \downarrow^ {u} \quad \downarrow^ {u ^ {\prime \prime}} \\ 0 \to \mathrm{N} ^ {\prime} \to \mathrm{N} \to \mathrm{N} ^ {\prime \prime} \to 0. \end{array}
\]

a commutative diagram of exact sequences, and let \(\tilde{u}':\mathbf{P}'\to\mathbf{Q}'\) and \(\tilde{u}'':\mathbf{P}''\to\mathbf{Q}''\) morphisms of complexes such that \(u'\circ p'=q'\circ\tilde{u}'\) and \(u''\circ p''=q''\circ\tilde{u}'\) . Show that there exists a morphism \(\tilde{u}:\mathbf{P}\to\mathbf{Q}\) such that \(u\circ p=q\circ\tilde{u}\) and that the diagram:

\[
0 \rightarrow \mathrm{P} ^ {\prime} \rightarrow \mathrm{P} \rightarrow \mathrm{P} ^ {\prime \prime} \rightarrow 0\tag{*}
\]

\[
0 \rightarrow Q ^ {\prime} \rightarrow Q \rightarrow Q ^ {\prime \prime} \rightarrow 0
\]

is commutative.

\begin{enumerate}
\def\labelenumi{\alph{enumi})}
\setcounter{enumi}{2}
\tightlist
\item
  Let \(\tilde{u}_{1}^{\prime}: \mathrm{P}^{\prime} \to \mathrm{Q}^{\prime}, \tilde{u}_{1}: \mathrm{P} \to \mathrm{Q}\) and \(\tilde{u}_{1}^{\prime\prime}: \mathrm{P}^{\prime\prime} \to \mathrm{Q}^{\prime\prime}\) three morphisms of complexes such that
\end{enumerate}

\[
u ^ {\prime} \circ p ^ {\prime} = q ^ {\prime} \circ \tilde {u} _ {1} ^ {\prime}, u \circ p = q \circ \tilde {u} _ {1}, u ^ {\prime \prime} \circ p ^ {\prime \prime} = q ^ {\prime \prime} \circ \tilde {u} _ {1} ^ {\prime \prime}
\]

making the diagram (*) commutative; let \(s'\) (resp. \(s''\) ) a homotopy connecting \(\tilde{u}'\) to \(\tilde{u}_{1}'\) (resp. \(\tilde{u}''\) to \(\tilde{u}_{1}'\) ). Show that there exists a homotopy \(s\) connecting \(\tilde{u}\) to \(\tilde{u}_{1}\) , such that the diagram:

\[
\begin{array}{c}0 \rightarrow \mathrm{P} ^ {\prime} \rightarrow \mathrm{P} \rightarrow \mathrm{P} ^ {\prime \prime} \rightarrow 0\\\downarrow s ^ {\prime} \quad \downarrow s \quad \downarrow s ^ {\prime \prime}\\0 \rightarrow \mathrm{Q} ^ {\prime} \rightarrow \mathrm{Q} \rightarrow \mathrm{Q} ^ {\prime \prime} \rightarrow 0\end{array}
\]

is commutative.

\begin{enumerate}
\def\labelenumi{\alph{enumi})}
\setcounter{enumi}{3}
\tightlist
\item
  State and prove the results corresponding to a), b), c) for injective resolutions. ¶17) In this exercise, if \((\mathbf{K}, d', d'')\) is a bicomplex (X, p. 174, exercise 11), we will denote \(K_{p,*}\) the complex \((\bigoplus_{n} K_{p,n}, d'')\) Every complex will be considered as a bicomplex, zero in degree \((p, q)\) for \(q \neq 0\) .
\end{enumerate}

Let C be a complex of A-modules. A projective resolution of C is a pair \((P, p)\) , where P is a bicomplex and \(p : P \to C\) a morphism of bicomplexes, such that \(P_{p,q} = 0\) for q \textless{} 0 and such that the complex \(P_{p,*}\) (resp. \(Z'_{p,*}(\hat{P})\) , \(B'_{p,*}(\mathbf{P})\) , \(H'_{p,*}(\mathbf{P})\) ) define a projective resolution of \(C_p\) (resp. \(Z_p(\mathbf{C})\) , \(B_p(\mathbf{C})\) , \(H_p(\mathbf{C})\) ) for every \(p \in Z\) . For \((P, p)\) be a projective resolution of C, it suffices that the complexes \(B'_{p,*}(\mathbf{P})\) and \(H'_{p,*}(\mathbf{P})\) define, for every p, projective resolutions of \(B_p(\mathbf{C})\) and \(H_p(\mathbf{C})\) .

\begin{enumerate}
\def\labelenumi{\alph{enumi})}
\tightlist
\item
  Show that every A-complex C admits a projective resolution (choose, for every p, projective resolutions \(B_{p,*}\) and \(H_{p,*}\) of \(B_{p}(C)\) and \(H_{p}(C)\) ; using Exercise 16, construct projective resolutions \(Z_{p,*}\) of \(Z_{p}(C)\) and \(P_{p,*}\) of \(C_{p}\) as well as exact sequences:
\end{enumerate}

\[
0 \rightarrow \mathrm{B} _ {p, *} \rightarrow \mathrm{Z} _ {p, *} \rightarrow \mathrm{H} _ {p, *} \rightarrow 0; \quad 0 \rightarrow \mathrm{Z} _ {p, *} \rightarrow \mathrm{P} _ {p, *} \rightarrow \mathrm{B} _ {p + 1, *} \rightarrow 0).
\]

\begin{enumerate}
\def\labelenumi{\alph{enumi})}
\setcounter{enumi}{1}
\item
  Let \(\mathbf{C}'\) another complex, \((\mathbf{P}', p')\) a projective resolution of \(\mathbf{C}'\) , \(u: \mathbf{C} \to \mathbf{C}'\) a morphism of complexes. Show that there exists a morphism of bicomplexes \(\tilde{u}: \mathbf{P} \to \mathbf{P}'\) such that \(u \circ p = p' \circ \tilde{u}\) (for every \(p \in \mathbf{Z}\) , choose morphisms \(\mathbf{B}_{p'}', (\mathbf{P}) \to \mathbf{B}_{p'}', (\mathbf{P}')\) and \(\mathrm{H}_{p'}', (\mathbf{P}) \to \mathrm{H}_{p'}', (\mathbf{P}')\) ; construct \(\tilde{u}\) as before, using Exercise 16, b)).
\item
  Let \(v: \mathbf{C} \to \mathbf{C}'\) a morphism of complexes homotopic to \(u\) , and let \(\tilde{v}: \mathbf{P} \to \mathbf{P}'\) be a morphism of bicomplexes such that \(v \circ p = p' \circ \tilde{v}\) . Show that \(\tilde{u}\) and \(\tilde{v}\) are homotopic (in the sense of X, p. 174, Exercise 11, c); use Exercise 16, c)).
\end{enumerate}

\(d)\) Let E be a subset of \(\mathbf{Z}\) , and C an A-complex such that \(\mathbf{C}_{p}=0\) for \(p\in\mathbf{E}\) . Show that there exists a projective resolution \((\mathbf{P},p)\) of C such that \(\mathbf{P}_{p,q}=0\) for \(p\in\mathbf{E}\) .

\begin{enumerate}
\def\labelenumi{\alph{enumi})}
\setcounter{enumi}{4}
\item
  Suppose there exists an integer \(n\) such that every A-module admits a projective resolution of length \(\leqslant n\) . Show that every A-complex C admits a projective resolution \((P, p)\) such that \(P_{p,q} = 0\) for \(q > n\) .
\item
  Let C be an A-complex. An injective resolution of C is a pair \((e, l)\) , where I is a bicomplex and \(e: C \to 1\) a morphism of bicomplexes, such that \(I^{p,q} = 0\) for q \textless{} 0 and such that the complex \(I^{p,*}\) (resp. \(Z^{p,*}(I)\) , \(B^{p,*}(l)\) , \(H^{p,*}(l)\) ) defines an injective resolution of \(C^{p}\) (resp. \(Z^{p}(C)\) , \(B^{p}(C)\) , \(H^{p}(C)\) ) for every \(p \in Z\) . State and prove the analogues of results a) to e) for injective resolutions.
\end{enumerate}

\begin{enumerate}
\def\labelenumi{\arabic{enumi})}
\setcounter{enumi}{17}
\tightlist
\item
  Assume the ring A is commutative. A differential graded A-algebra is an A-algebra S, graded of type Z, such that \(S_{n} = 0\) for n \textless{} 0, equipped with an antiderivation d of degree (-1), satisfying \(d \circ d = 0\) . We shall also denote by S the complex of A-modules (S, d).
\end{enumerate}

\begin{enumerate}
\def\labelenumi{\alph{enumi})}
\item
  Let C be a complex of A-modules, zero in degrees \textless{} 0. Show that there exists on S = T(C) (resp. S = S(C), resp. S = Λ(C)) a structure of differential graded A-algebra, such that the canonical injection of C into S is a morphism of complexes.
\item
  Show that the multiplication of S induces a structure of graded A-algebra on H(S).
\item
  Let S, T be two differential graded A-algebras. Set D(s ⊗ t) = ds ⊗ t + (-1)\^{}p s ⊗ dt for s ∈ S\_p, t ∈ T. Show that D defines a structure of differential graded A-algebra on the left tensor product S *⊗\_A T (III, p. 49).
\end{enumerate}

\begin{enumerate}
\def\labelenumi{\arabic{enumi})}
\setcounter{enumi}{18}
\tightlist
\item
  Assume the ring A is commutative.
\end{enumerate}

\begin{enumerate}
\def\labelenumi{\alph{enumi})}
\item
  Let \((\mathbb{S}, d_{\mathbb{S}})\) a differential graded A-algebra (exercise 18) , M an A-module, \(u: \mathbb{M} \to Z_{n}(\mathbb{S})\) an A-homomorphism. Show that there exists on \(\mathbb{S} \otimes_{\mathbb{A}} \mathbf{T}(\mathbb{M})\) a structure of differential graded A-algebra such that \(d(s \otimes 1) = d_{\mathbb{S}} s \otimes 1\) for \(s \in \mathbb{S}\) and \(d(1 \otimes m) = u(m) \otimes 1\) for \(m \in \mathbb{M}\) .
\item
  Let B be an A-algebra. Show that there exists a differential graded A-algebra S and a homomorphism of A-algebras \(p: \mathbf{S}_{0} \to \mathbf{B}\) such that \((\mathbf{S}, p)\) be a free resolution of the A-module B. (Construct by induction on \(n\) , using \(a\) ), a differential graded algebra \(\mathbf{S}(n)\) , free in each degree, such that \(\mathrm{H}_{i}(\mathbf{S}(n)) = 0\) for \(0 < i < n\) and \(\mathrm{H}_{0}(\mathbf{S}(n)) = \mathbf{B}\) .)
\item
  If the algebra B is commutative, show that one can find an alternating A-algebra S (III, p. 53) satisfying the conditions of b). If moreover A is noetherian and if B is a quotient of A by an ideal, show that one can find S such that \(S_{0} = A\) and that \(S_{n}\) is a finitely generated A-module for every n.
\end{enumerate}

¶ 20) Assume that the ring A is an algebra over a commutative ring k. Let \(\eta: k \to A\) be the homomorphism defined by \(\eta(\lambda) = \lambda 1_{A}\) for \(\lambda \in k\) , and let \(\overline{A}\) its cokernel. If \(a\) is an element of A, we denote by \(\overline{a}\) its class in \(\overline{A}\) . We set \(\mathrm{N}_{n}(\mathrm{A}) = \mathrm{A} \otimes_{k} \overline{\mathrm{A}}^{\otimes n} \otimes_{k} \mathrm{A}\) for \(n \geqslant 0\) , \(\mathrm{N}_{n}(\mathrm{A}) = 0\) for \(n < 0\) , and

\[
\mathrm{N} (\mathbf {A}) = \bigoplus_ {n \in \mathbf {Z}} \mathrm{N} _ {n} (\mathbf {A}).
\]

\begin{enumerate}
\def\labelenumi{\alph{enumi})}
\tightlist
\item
  Show that there exists an (A, A)-endomorphism d of N(A), graded of degree (-1), such that
\end{enumerate}

\[
\begin{array}{l} d (a \otimes \bar {a} _ {1} \otimes \dots \otimes \bar {a} _ {n} \otimes b) = a a _ {1} \otimes \bar {a} _ {2} \otimes \dots \otimes \bar {a} _ {n} \otimes b + \sum_ {i = 1} ^ {n - 1} (- 1) ^ {i} a \otimes \dots \otimes \overline {{a _ {i}}} \overline {{a _ {i + 1}}} \otimes \dots \otimes \bar {a} _ {n} \otimes b \\ \qquad + (- 1) ^ {n} a \otimes \bar {a} _ {1} \otimes \dots \otimes \bar {a} _ {n - 1} \otimes a _ {n} b \quad \text {for} \quad a, a _ {1},..., a _ {n}, b \in A \end{array}
\]

\begin{enumerate}
\def\labelenumi{\alph{enumi})}
\setcounter{enumi}{1}
\item
  Show that \(d \circ d = 0\) , so that \(\mathbf{N}(\mathbf{A})\) is a complex of \((\mathbf{A}, \mathbf{A})\) -bimodules. If M is a left A-module, we denote by \(\mathbf{N}(\mathbf{A}, \mathbf{M})\) the A-complex \(\mathbf{N}(\mathbf{A}) \otimes_{\mathbf{A}} \mathbf{M}\) .
\item
  Let B(A, M) be the standard resolution of M (X, p. 58), \(p: \mathrm{B}(\mathrm{A}, \mathrm{M}) \to \mathrm{N}(\mathrm{A}, \mathrm{M})\) the A-homomorphism defined by \(p(a_0 \otimes \cdots \otimes a_n \otimes m) = a_0 \otimes \overline{a}_1 \otimes \cdots \otimes \overline{a}_n \otimes m\) for \(a_0, \ldots, a_n \in \mathbf{A}, m \in \mathbf{M}\) . Show that \(p\) is a homotopy equivalence of A-complexes (construct, by induction, a map inverse up to homotopy).
\end{enumerate}

* 21) Let \(k\) be a commutative ring, and let g be a \(k\) Lie algebra, U its enveloping algebra (cf. LIE, I, § 2). Assume that g is a \(k\) -free module.

\begin{enumerate}
\def\labelenumi{\alph{enumi})}
\tightlist
\item
  Show that for every \(n \geqslant 0\) an injective U-homomorphism
\end{enumerate}

\[
j _ {n}: \mathrm{U} \otimes_ {k} \Lambda^ {n} (\mathrm{g}) \rightarrow \mathrm{U} ^ {\otimes (n + 1)} \quad \text { such that } \quad j _ {n} (u \otimes (x _ {1} \wedge \dots \wedge x _ {n})) = \sum_ {\sigma \in \mathfrak {S} _ {n}} \varepsilon (\sigma), u \otimes x _ {\sigma (1)} \otimes \dots \otimes x _ {\sigma (n)}
\]

for \(u\in \mathbf{U},x_1,\dots ,x_n\in \mathfrak{g}\)

We thus define a graded U-homomorphism of degree 0. \(j:U\otimes_{k}\Lambda(g)\to B(U,k)\) , where \(B(U,k)\) is the standard resolution of the left U-module \(k\) (the structure of U-module of \(k\) being deduced from the natural homomorphism \(U\to k\) ). Show that \(U\otimes_{k}\Lambda(g)\) is identified by \(j\) with a subcomplex of \(B(U,k)\) , the differential being given by the formula :

\[
\begin{array}{l} d (u \otimes (x _ {1} \wedge \dots \wedge x _ {n})) = \sum_ {1 \leqslant i \leqslant n} (- 1) ^ {i + 1} u x _ {i} \otimes (x _ {1} \wedge \dots \wedge \hat {x} _ {i} \wedge \dots \wedge x _ {n}) \\ + \sum_ {1 \leqslant i <   j \leqslant n} (- 1) ^ {i + j} u \otimes ([ x _ {i}, x _ {j} ] \wedge x _ {1} \wedge \dots \wedge \hat {x} _ {i} \wedge \dots \wedge \hat {x} _ {j} \wedge \dots \wedge x _ {n}) \quad \text { pour } \quad u \in U, x _ {1},..., x _ {n} \in g. \end{array}
\]

(where the sign over a letter means that it is to be omitted). We denote by V(g) the U-complex thus defined. b) Show that V(g) defines a free resolution of the U-module k (let \((\mathrm{U}_{n})_{n\geqslant0}\) be the natural filtration of U, \(F_{r}V(g)\) the sub-k-module of V(g) generated by the elements \(u_{p}\otimes x_{q}\) for \(u_{p}\in U_{p}, x_{q}\in\Lambda^{q}(g), p+q\leqslant r\) . Show that \(F_{r}V(g)\) is a subcomplex of V(g), and that the complex \(\bigoplus_{r\geqslant0}F_{r}V(g)/F_{r-1}V(g)\)

is isomorphic to the complex \(\mathbf{S}(\mathfrak{g})\otimes\symbf{\Lambda}(\mathfrak{g})\) defined in X, p. 151). *
